\documentclass[11pt,a4paper]{article}

% ── Geometry ──────────────────────────────────────────────────────────────────
\usepackage[margin=1in,headheight=14pt]{geometry}

% ── Core packages ─────────────────────────────────────────────────────────────
\usepackage{amsmath,amssymb,amsthm}
\usepackage{mathrsfs}
\usepackage{graphicx}
\usepackage{booktabs}
\usepackage{xcolor}
\usepackage{enumitem}
\usepackage{caption}
\usepackage{subcaption}
\usepackage{needspace}
\usepackage{titlesec}
\usepackage{fancyhdr}
\usepackage{doi}
\usepackage[expansion=false]{microtype}

% ── natbib ────────────────────────────────────────────────────────────────────
\usepackage{natbib}
\setcitestyle{numbers,square,sort&compress}

% ── hyperref (must be near-last) ──────────────────────────────────────────────
\usepackage{hyperref}
\hypersetup{
  colorlinks=true,
  linkcolor=cyan!70!black,
  citecolor=cyan!70!black,
  urlcolor=cyan!70!black,
  pdftitle={TMRB Soft-Hair Resolution of the Black-Hole Information Paradox},
  pdfauthor={George Bressler},
}

% ── Colours ───────────────────────────────────────────────────────────────────
\definecolor{navyblue}{HTML}{0D1B2A}
\definecolor{gold}{HTML}{C8A84B}
\definecolor{cyancol}{HTML}{00B4D8}

% ── Section formatting ────────────────────────────────────────────────────────
\titleformat{\section}
  {\needspace{6\baselineskip}\large\bfseries\color{navyblue}}
  {\thesection}{1em}{}[\color{gold}\titlerule]
\titleformat{\subsection}
  {\needspace{4\baselineskip}\normalsize\bfseries\color{navyblue}}
  {\thesubsection}{1em}{}

% ── Abstract fix (prevent margin bleed) ───────────────────────────────────────
\renewenvironment{abstract}{%
  \begin{quote}\small\textbf{Abstract.}\quad\ignorespaces}%
  {\end{quote}}

% ── Typography ────────────────────────────────────────────────────────────────
\setlength{\emergencystretch}{2em}

% ── Header / Footer ───────────────────────────────────────────────────────────
\pagestyle{fancy}
\fancyhf{}
\fancyhead[L]{\small\color{navyblue}\textbf{BRESSLER 2026 · PTRH PAPER 10 · TMRB SOFT-HAIR BH INFORMATION · ZENODO PREPRINT · CC-BY 4.0}}
\fancyfoot[C]{\small Page \thepage\ · doi:10.5281/zenodo.19981840 · Open access CC-BY 4.0}
\renewcommand{\headrulewidth}{0.4pt}
\renewcommand{\footrulewidth}{0.4pt}

% ── Theorem env ───────────────────────────────────────────────────────────────
\newtheorem{theorem}{Theorem}
\newtheorem{proposition}[theorem]{Proposition}
\newtheorem{conjecture}[theorem]{Conjecture}

% ─────────────────────────────────────────────────────────────────────────────
\begin{document}

% ── Title block ───────────────────────────────────────────────────────────────
\title{%
  \textbf{TMRB Soft-Hair Resolution of the Black-Hole Information Paradox:}\\[4pt]
  \textbf{Discrete $\mathbb{Z}_9$ Torsion Hair, Finite Microstate Counting,}\\[4pt]
  \textbf{and the $9\ln 9$ Topological Correction in 11D M-Theory}%
}
\author{%
  George Bressler\\[2pt]
  ORCID: \href{https://orcid.org/0009-0001-9439-4930}{0009-0001-9439-4930}\\
  \href{mailto:george.bressler@proton.me}{\texttt{george.bressler@proton.me}}\\
  @GeorgeBressler · Sarasota, FL
}
\date{April 2026 (v2: May 2026; v3: May 2026; v4: November 2026)}
\maketitle
\thispagestyle{fancy}

% ── Prologue ──────────────────────────────────────────────────────────────────
\vspace{4pt}
\begin{center}
{\color{gold}\rule{0.38\textwidth}{0.4pt}}
\end{center}
\vspace{2pt}

\begin{quote}
\itshape
We think of black holes as the universe's ultimate shredders. You fall in,
you are gone, your information destroyed, end of story. Except the laws of
quantum mechanics say information cannot be destroyed --- not as a suggestion,
but as a requirement as fundamental as anything physics has established. So
we have a paradox. The most powerful objects in the universe appear to
violate its most inviolable rule.

Here is where it gets beautiful.

Imagine the surface of a black hole --- the event horizon --- as the world's
most sophisticated recording medium. When you press a vinyl record, the stylus
cuts grooves into the surface: information permanently encoded in the physical
structure of the disc. Every infalling particle, every photon, every quantum of
information that crosses the horizon does the same thing. It cuts a groove. A
gravitational groove. A permanent deformation of the quantum vacuum at the
surface of spacetime itself. Physicists call these imprints \emph{soft hair} ---
because the horizon appears smooth and featureless from the outside, just as
Hawking's equations predict, but at the quantum level it is covered in an
impossibly fine structure that remembers everything that ever fell in.

The question that went unanswered for fifty years is whether that record can
actually be \emph{read}. Hawking's calculation says the radiation emerging from
a black hole is pure thermal noise --- a cosmic hiss carrying no signal, no
memory, no message. The grooves are there, but the record is unplayable.
Information encoded but permanently inaccessible, which is almost worse than
destroyed.

This paper proposes an answer to that question. The record is not analog. The grooves do not
form a continuous spectrum. The torsion structure of the geometry underlying
the compactification of eleven-dimensional spacetime divides the entire
recording surface into exactly nine orthogonal channels per torsion cell
(about nine Planck areas) of the horizon. Every possible quantum state that falls into a black hole is
encoded as a chord built from those nine fundamental modes. The radiation coming
out is not noise --- it is the playback, one quantum at a time, through those
same nine channels, in the order the information arrived --- that is the
claim this paper proposes. The
$\mathbb{Z}_9$ conservation law is the bookkeeping every playback must respect;
showing that it is enough to lose no note is the open task (Gap~11).

If that task is done, the universe kept the receipt.

All of this follows from a single observational anchor: one measurement of the
Milky Way's phase spiral by the Gaia satellite, one number,
$\eta_{\mathrm{diff}} = 0.82 \pm 0.02$, which fixes the framework's
normalisation and connects the largest observable structure in our galaxy to the
smallest meaningful scale in physics. Whether that link is more than a
coincidence is what the predictions below are built to test.

What follows is the derivation.
\end{quote}

\vspace{2pt}
\begin{center}
{\color{gold}\rule{0.38\textwidth}{0.4pt}}
\end{center}
\vspace{6pt}

% ── Abstract ──────────────────────────────────────────────────────────────────
\begin{abstract}
The TMRB (Torsion-Modulated Recursive Branching) framework—built on 11D
M-theory compactified on the Schoen gyroid with $\mathbb{Z}_9$ discrete
torsion and modular-covariant Gaussian branching weights
$w_n = C\exp(-\pi d(n)^2/9)$—has proposed a discrete holographic dual
for flat spacetime and an explicit closure of the infrared triangle linking the
Weinberg soft theorem, BMS supertranslation symmetry, and gravitational memory.
This work proposes a resolution of the black-hole information paradox within
the identical construction, and states which of its steps remain open.

The event horizon carries exactly nine quantised soft-hair charge sectors per
torsion cell from the $\mathbb{Z}_9$ gyroid torsion nodes, yielding a finite
horizon Hilbert space of dimension
$\dim\mathcal{H}_{\mathrm{hair}} = 9^{N_{\mathrm{cell}}} = e^{S_{\mathrm{BH}}}$,
where $N_{\mathrm{cell}} = A/(4G_4\ln 9)$ torsion cells tile the horizon.
Beyond the one-loop logarithm $c_{\log}\ln(A/G_4)$, the Bekenstein--Hawking
entropy receives a size-independent topological correction
\[
  S_{\mathrm{TMRB}} = \frac{A}{4G_4} + c_{\log}\ln\frac{A}{G_4} + 9\ln 9 + \ln 3
    + \mathcal{O}\!\left(\frac{G_4}{A}\right),
  \qquad \Delta S_{\mathrm{tors}} = 9\ln 9 + \ln 3 \approx 20.874,
\]
where $9\ln 9 \approx 19.775$ is the nine-node torsion hair-lattice log-degeneracy
and $\ln 3 \approx 1.099$ is the Euclidean boundary contribution from
$\ln Z_{\rm tors}$.  A golden-ratio factor
$\varphi^{-N_*}$ ($\varphi = (1+\sqrt{5})/2$) is proposed to suppress
logarithmic soft-graviton loop corrections; it is not yet derived from the
weights (Section~\ref{sec:loop}, Gap~15), while
holographic consistency is argued from torus partition function factorisation
$Z(\tau)=|\chi_{\mathbb{Z}_9}(\tau)|^2$ (Section~\ref{sec:unitarity}).
Information-carrying remnants are excluded by a capacity bound of
$\Delta S_{\rm tors} \approx 20.9$ nats; whether the construction
avoids firewalls, and the full unitary map from infalling state to radiation,
are open (Section~\ref{sec:gaps}, Gaps~10 and~11).  All results anchor to the single Gaia DR3
calibration $\eta_{\mathrm{diff}} = 0.82\pm 0.02$.  The 11D CDT extension
invoked here is conjectured; the framework is non-supersymmetric.  Four
correlated predictions are presented (three near-term falsifiable; the entropy gap
has no observable yet), together with
two sensitivity scales: GRACE/GRACE-FO mascon gravity (about
$4\times10^{-9}~\mathrm{m\,s^{-2}}$, roughly $10^{6}$ times above the framework's predicted
quasi-static bulk-gravity-leakage amplitude), and a GWOSC O3a
narrowband sensitivity floor that lies below that amplitude but applies only to a
component near 100~Hz, which the prediction does not contain. The ringdown
prediction was tested against published LIGO--Virgo--KAGRA data by the X24si method
(registered with an outside clock before any result was read, and run once): GWTC-5.0
gives $\delta\hat\tau_{220} = 0.07^{+0.06}_{-0.05}$, which excludes the predicted
shortening, so under its stated assumptions that prediction is disfavoured.

\smallskip
\noindent\textbf{Keywords:} black-hole information paradox, soft hair, BMS symmetry,
torsion, M-theory, CDT, gyroid, holography, Hawking radiation, unitarity.
\end{abstract}

% ── AI Disclosure ─────────────────────────────────────────────────────────────
{\small\textit{AI disclosure: Mathematical iteration via Grok (xAI);
document formatting and refinement via Claude (Anthropic).
All theoretical content, interpretations, and responsibility remain with the author.}}

\vspace{8pt}

% ─────────────────────────────────────────────────────────────────────────────
\needspace{6\baselineskip}
\section{Introduction}
\label{sec:intro}

The black-hole information paradox~\cite{Hawking1975} has resisted resolution for
five decades.  At its core lies a conflict: Hawking's semiclassical calculation
predicts thermal radiation and entropy loss, while unitarity of quantum mechanics
demands that information escaping a black hole remains accessible in principle.
Proposed resolutions range from Planck-scale remnants~\cite{AharonovEtal1987} and
black-hole complementarity~\cite{Susskind1993} to the firewall
conjecture~\cite{AMPS2013} and island-formula approaches~\cite{Penington2020}.

The soft-hair programme of Hawking, Perry, and Strominger
(HPS)~\cite{HawkingPerryStrominger2016} proposed that the horizon carries an
infinite-dimensional memory of soft graviton and soft photon charges, encoding
infalling information.  While compelling, the original proposal left open how the
soft-hair Hilbert space is regularised and how it connects to microscopic string or
quantum-gravity degrees of freedom.

The TMRB (Torsion-Modulated Recursive Branching) framework~\cite{BresslerUF2026,BresslerPrismatic2026}
provides precisely this regularisation.  Compactifying 11D M-theory on the Schoen
gyroid manifold~\cite{Schoen1970} with $\mathbb{Z}_9$ discrete torsion~\cite{BresslerBGL2026}
discretises the Bondi--Metzner--Sachs (BMS) superrotation group to nine generators,
closing the infrared
triangle~\cite{BresslerCH2026} via Weinberg's soft theorem~\cite{Weinberg1965},
BMS supertranslations~\cite{BMS1962,Strominger2014}, and gravitational
memory~\cite{Christodoulou1991,BresslerGP2026}.  The same 3-6-9 torsion
mathematics that sets the 230~Myr galactic disk-wave period and the $\mathbb{Z}_9$
parafermion spectrum at null infinity~\cite{GepnerQiu1987,ZamoFateev1987} now
threads the black-hole event horizon.

This paper argues that:
\begin{enumerate}[leftmargin=*, label=\arabic*.]
  \item The horizon carries exactly \emph{nine} quantised soft-hair charge sectors
    per torsion cell from the $\mathbb{Z}_9$ torsion nodes (Section~\ref{sec:hair}).
  \item The Bekenstein--Hawking entropy acquires, beyond the one-loop logarithm,
    a universal size-independent torsion correction $\Delta S_{\mathrm{tors}} = 9\ln 9 + \ln 3 \approx 20.874$
    (Section~\ref{sec:entropy}).
  \item A golden-ratio factor $\varphi^{-N_*}$ is proposed to suppress logarithmic
    soft-graviton loop corrections (Section~\ref{sec:loop}); it is not yet derived
    from the weights $w_n = C\exp(-\pi d(n)^2/9)$ (Gap~15).
  \item The torus partition function factorises in the discrete holomorphic sector,
    supporting holographic consistency (Section~\ref{sec:unitarity}; Gap~13).
\end{enumerate}

Open theoretical gaps are documented in Section~\ref{sec:gaps}.
Four predictions (three falsifiable near-term) and current observational constraints are
presented in Section~\ref{sec:predictions}.
For the complete series cross-citation map, master parameter table, and
TRL-3/4 test protocol spanning all eight active papers, see~\cite{BresslerSurvey2026}.

\smallskip
\noindent\textbf{Changes in v2:}
\textit{(i)}~BMS spelled out as Bondi--Metzner--Sachs (BMS) on first body use.
\textit{(ii)}~CDT spelled out as Causal Dynamical Triangulations (CDT) on first body use, with explicit AJL citations at definition point.
\textit{(iii)}~$\Delta\lambda$ corrected to 1.7287 (from 1.729) throughout, consistent with all companion papers.
\textit{(iv)}~$Z_{\rm tors} \approx 2$ (not exactly~2; numerical deviation ${\approx}3\times10^{-4}$) corrected at two points.
\textit{(v)}~All cross-citation DOIs updated to canonical versions: TMRB Unified Framework $\to$ v3 (20248650), Celestial Holography $\to$ v4 (20249096), M-Theory Bulk Gravity Leakage $\to$ v10a (20276554), Galactic Piano $\to$ v4 (20262175), QGGPf Predictions $\to$ v4 (20272654).
\textit{(vi)}~Paper~1 (prismatic curvature foundation,~\href{https://doi.org/10.5281/zenodo.20146627}{20146627}) and Paper~9 (PTRH survey,~\href{https://doi.org/10.5281/zenodo.20315669}{20315669}) added to cross-citations.
\textit{(vii)}~\texttt{Vafa1994} bibitem key corrected to \texttt{Vafa1986} (paper is Nucl.~Phys.~B 1986).
\textit{(viii)}~NTRS URL added to Schoen (1970) reference.
\textit{(ix)}~Equation references standardised to \texttt{Eq.\textasciitilde\textbackslash eqref\{\}} throughout.
\textit{(x)}~Footer DOI set to 10.5281/zenodo.19981841.
\textit{(xi)}~\texttt{microtype[expansion=false]} added; \texttt{emergencystretch} corrected to 2em.
\textit{(xii)}~Gap~8 added: derivation of $G_4 = \ell_{\rm Pl}^2/\ln 9$ from the prismatic action is an open item (highest-priority tractable gap in PTRH programme).
\textit{(xiii)}~Gap~9 added: pre-Page-time retrieval protocol remains open.

\smallskip
\noindent\textbf{Changes in v3:}
\textit{(i)}~Horizon Hilbert space formula corrected for Bekenstein consistency:
$\dim\mathcal{H}_{\mathrm{hair}} = 9^{A/(4\ell_{\mathrm{Pl}}^2)}$ (v1/v2) replaced by
$\dim\mathcal{H}_{\mathrm{hair}} = 9^{N_{\mathrm{cell}}} = e^{S_{\mathrm{BH}}}$,
$N_{\mathrm{cell}} = A/(4G_4\ln 9)$ throughout (abstract, \S\ref{sec:hair},
conclusions); the v1/v2 form exceeds the Bekenstein--Hawking bound by a factor
of $\ln 9 \approx 2.197$ unless the identification $G_4 = \ell_{\mathrm{Pl}}^2/\ln 9$
(Gap~8) is assumed, which is noted explicitly.
\textit{(ii)}~Dangling cross-reference \texttt{app:microstate} removed; the
relevant argument is now inline at Eq.~\eqref{eq:N_hair}.
\textit{(iii)}~All cross-citation DOIs updated to concept DOIs (always resolve
to latest version) throughout.

\smallskip
\noindent\textbf{Changes in v4 (from v3, the last version on Zenodo):}
\begin{enumerate}[leftmargin=*, label=\textit{(\roman*)}, itemsep=1pt, topsep=2pt]
  \item \textbf{Ringdown prediction tested against published data and disfavoured, by
    the X24si method} (human-in-the-loop self-recursive reasoning, anchored by an outside
    clock and outside data). Before any result was read, a test of this paper's ringdown
    prediction against LIGO--Virgo--KAGRA data was registered with an outside clock (a
    timestamped GitHub commit): the claim, the data source, the decision rule and a frozen
    decision script. Run once on GWTC-5.0, the combined result
    $\delta\hat\tau_{220} = 0.07^{+0.06}_{-0.05}$ excludes the predicted shortening
    ($-0.095$; band $-0.126$ to $-0.070$), and all five secondary combinations agree. The
    result stands as run (Section~\ref{sec:x24si_test}).
  \item \textbf{Entropy correction revised.} The weights are made $\mathbb{Z}_9$-symmetric,
    so $\ln Z_{\rm tors} = \ln 3$ and $\Delta S_{\rm tors} = 9\ln 9 + \ln 3 \approx 20.874$
    (v3: $\ln 2$, 20.468). The one-loop logarithm $c_{\log}\ln(A/G_4)$, with
    $c_{\log} = 77/90$, is added; it exceeds $\Delta S_{\rm tors}$ unless Gap~15 closes.
  \item \textbf{Claims matched to what is shown.} ``Resolves'' becomes ``proposes a
    resolution''. Firewall avoidance, the unitary map, orbit counting, the 45-primary
    partition function, the remnant endpoint, the golden-ratio loop factor, the entropy-gap
    observable and the ringdown assumptions are listed as open Gaps~10--18. Only
    information-carrying remnants are excluded.
  \item \textbf{Ringdown prediction restated.} The nine sidebands are unresolved (0.094
    line widths at every mass); the prediction is a shortening of the damping time, given
    by remnant spin and fit start time in Table~\ref{tab:ringdown_spin} (exact Kerr
    frequencies).
  \item \textbf{Observational sections rewritten from the analysis scripts.} The GWOSC
    result is a sensitivity floor ($1.6\times10^{-16}~\mathrm{m\,s^{-2}}$, for a
    ${\sim}100$~Hz component only). The GRACE run of 2026-04-26 is described as done: its
    zero result follows from its method, and $4\times10^{-9}~\mathrm{m\,s^{-2}}$ is the
    1~cm water scale. A synthetic histogram is removed from Fig.~5; the run's own output
    is Fig.~6.
  \item \textbf{Corrections.} Torsion cell area $8.79\,\ell_{\rm Pl}^2$; sideband weights;
    loop bound; torsion classes in $H^3$ (Dijkgraaf--Witten); modular invariance as a
    consistency condition; imaginary $\tau$; quantum dimension 3; charges modulo $2\pi$;
    footer DOI set to the concept DOI; references corrected and added.
  \item The full revision trail (r1--r10), audits, registration files and analysis scripts
    are included in the Zenodo record.
\end{enumerate}

% ─────────────────────────────────────────────────────────────────────────────
\needspace{6\baselineskip}
\section{TMRB Framework Recap}
\label{sec:recap}

\subsection{Canonical Parameters}

All TMRB results are anchored to a single calibration from Gaia DR3~\cite{GaiaDR3}:
\begin{equation}
  \eta_{\mathrm{diff}} = 0.82 \pm 0.02.
  \label{eq:etadiff}
\end{equation}
The remaining parameters are derived or independently motivated (Table~\ref{tab:params}).

\begin{table}[ht]
\centering
\caption{TMRB canonical parameters. $\kappa_{\mathrm{CE}}$ and $\kappa_{\mathrm{prop}}$
are \emph{distinct} quantities; the former governs conformal-weight combs and
galactoseismology, the latter governs plasma torsion suppression and IAW lock.}
\label{tab:params}
\begin{tabular}{llp{4.0cm}l}
\toprule
Symbol & Value & Role & Domain \\
\midrule
$\eta_{\mathrm{diff}}$ & $0.82\pm0.02$ & Phase-space diffusion ratio & All papers \\
$\lambda_{\mathrm{rec}}$ & 5.2 & Recursive branching depth & All papers \\
$\mu$ & 2.1 & Branching modulation & All papers \\
$\kappa_{\mathrm{CE}}$ & $\mu/\lambda_{\mathrm{rec}} = 0.4038$ & Conformal-weight comb & Holography, galactoseismology \\
$\kappa_{\mathrm{prop}}$ & 0.48 & Torsion suppression exponent & Plasma sector, QGGPf \\
$d_{\mathrm{opt}}$ & 9 & Self-tuned branching depth & All papers \\
$\varphi$ & $(1+\sqrt{5})/2$ & Golden ratio & Branching weights \\
$G_4/G_{11}$ & $(5.214\pm0.12)\times10^{-40}$ & KK hierarchy ratio (dimensionful: $G_{11}$ carries length$^9$; units as in~\cite{BresslerBGL2026}) & Holography, QGGPf \\
$T_0$ & 230~Myr & Galactic disk-wave crest period & Galactoseismology \\
$\eta_{\mathrm{null}}$ & $1.0\pm0.15$ & Null model baseline & All falsification tests \\
\bottomrule
\end{tabular}
\end{table}

\subsection{Gyroid Compactification and $\mathbb{Z}_9$ Torsion}

M-theory is compactified on the Schoen gyroid, a triply-periodic minimal surface
with crystallographic space group $\Gamma = Ia\bar{3}d$ (No.~230)~\cite{Schoen1970}.
The gyroid has no G$_2$ holonomy; the construction is explicitly
non-supersymmetric, with UV finiteness arising from the causal structure of the
conjectured 11D Causal Dynamical Triangulations (CDT) extension~\cite{AJL2004,AJL2005}.

The discrete torsion phases are
\begin{equation}
  \theta_n = \frac{2\pi n}{9}, \qquad n = 0,1,\ldots,8,
  \label{eq:torsion_phases}
\end{equation}
labelled by the classes of $H^3(\mathbb{Z}_9, U(1)) \cong \mathbb{Z}_9$, the
Dijkgraaf--Witten classification of $\mathbb{Z}_9$ three-cocycles~\cite{DijkgraafWitten1990},
the natural home of a discrete C-field phase. (Vafa's discrete
torsion~\cite{Vafa1986} is classified by $H^2(G, U(1))$, which is trivial for cyclic
$G = \mathbb{Z}_9$; ``torsion'' in this paper means the $H^3$ class.)  The modular-covariant
branching weights are
\begin{equation}
  w_n = C\sum_{k\in\mathbb{Z}}\exp\!\left(-\frac{\pi (n+9k)^2}{9}\right)
      \simeq C\exp\!\left(-\frac{\pi\, d(n)^2}{9}\right),
  \qquad d(n) = \min(n,\,9-n),
  \label{eq:wn}
\end{equation}
the Gaussian periodised on $\mathbb{Z}_9$. It is an exact fixed point of the
$\mathbb{Z}_9$ discrete Fourier transform,
$w_n = \frac{1}{3}\sum_{m} w_m e^{-2\pi i mn/9}$ (checked numerically to $10^{-15}$),
and it respects $\mathbb{Z}_9$ conjugation, $w_{9-n} = w_n$. The one-sided form
$\exp(-\pi n^2/9)$, $n = 0,\dots,8$, used through v3 is neither: it misses the
fixed point by $0.57$ and gives $w_8 \ll w_1$. The periodised form ensures self-consistency under the modular
$S$-transformation $\tau\to-1/\tau$.  This property replaces the earlier
$\varphi^{-\lambda n}$ form, which lacked modular covariance.
The resulting torsion partition function and associated quantities are:
\begin{align}
  Z_{\mathrm{tors}} &= \sum_{n=0}^{8}\sum_{k\in\mathbb{Z}} e^{-\pi (n+9k)^2/9}
    = \sum_{m\in\mathbb{Z}} e^{-\pi m^2/9} = \vartheta_3\!\left(0\,\middle|\,\tfrac{i}{9}\right)
    = \frac{1}{C} = 3 + \mathcal{O}\!\left(e^{-9\pi}\right),
  \label{eq:Z_tors_val}\\
  \mathcal{S}_{\mathrm{tors}} &= -\sum_{n=0}^{8} w_n \ln w_n \approx 1.598,
  \label{eq:torsion_entropy}\\
  \mathcal{H}_{\mathrm{hair}} &\equiv \ln(9^9) = 9\ln 9 \approx 19.775.
  \label{eq:H_hair}
\end{align}
By Poisson summation, $\sum_{m\in\mathbb{Z}} e^{-\pi m^2/9} = 3\sum_{k\in\mathbb{Z}} e^{-9\pi k^2}$,
so $Z_{\rm tors} = 3$ to about $10^{-12}$; $3 = \sqrt{9}$ is also the total quantum
dimension of $\mathbb{Z}_9$ (Section~\ref{sec:entropy}). Through v3 the one-sided sum
was used, which gives 2 and $\ln 2$; it breaks $\mathbb{Z}_9$ conjugation and is withdrawn.
Here $\mathcal{S}_{\mathrm{tors}}$ is the Shannon entropy of the Gaussian weight
distribution; $\mathcal{H}_{\mathrm{hair}} = 9\ln 9$ is the log-degeneracy of
the full nine-node torsion hair lattice (the number of distinct $\mathbb{Z}_9$
configurations across all nine nodes), equal to $\ln(9^9)$; and
$\ln Z_{\mathrm{tors}} = \ln 3$ is the boundary contribution to the gravitational
path integral.  These three quantities are distinct and are not conflated in what
follows.

\subsection{TMRB Partition Function}

The full partition function is
\begin{equation}
  Z = \sum_{\{\Delta\}} \prod_{\mathrm{links}} e^{i\theta_n}
      \prod_{\mathrm{vertices}} A_{\mathrm{EPRL/FK}}
      \prod_{\mathrm{branches}} w_n,
  \label{eq:Z}
\end{equation}
where $A_{\mathrm{EPRL/FK}}$ are the EPRL/FK spin-foam vertex amplitudes and the
sum runs over all causal dynamical triangulations $\{\Delta\}$ consistent with the
causal structure (11D CDT extension conjectured).

% ─────────────────────────────────────────────────────────────────────────────
\needspace{6\baselineskip}
\section{$\mathbb{Z}_9$ Soft Hair on the Event Horizon}
\label{sec:hair}

\subsection{BMS Superrotations Discretised to $\mathbb{Z}_9$}

In the existing Celestial Holography paper~\cite{BresslerCH2026}, the nine torsion
nodes discretise the BMS superrotation group at null infinity $\mathscr{I}^+$ to
nine generators.  The same mechanism operates at the horizon, which shares the
BMS symmetry structure under Rindler expansion~\cite{Donnay2016}: the near-horizon
geometry in Rindler coordinates $(\rho, t, x^i)$ admits supertranslation charges
$T(z,\bar{z})$ and superrotation charges $Y^z(z)$ on the bifurcation 2-sphere.
The $\mathbb{Z}_9$ torsion phases $\theta_n$ inherited from the bulk
compactification label nine independent charge sectors on the bifurcation surface
by the following argument: the 11D C-field flux restricted to any 3-cycle
$\Sigma_3$ threading the horizon carries the phase $e^{i\theta_n}$; upon
dimensional reduction, this phase labels distinct boundary states on the
bifurcation $S^2$ exactly as at $\mathscr{I}^+$.  This identification follows
directly from the Kaluza--Klein reduction of the 11D C-field on the gyroid:
the discrete torsion classes in $H^3(\mathbb{Z}_9, U(1)) \cong \mathbb{Z}_9$
~\cite{DijkgraafWitten1990} are preserved by the reduction and project
onto the bifurcation $S^2$ as $\mathbb{Z}_9$-valued boundary phases, mirroring
the discretisation at $\mathscr{I}^+$ established in~\cite{BresslerCH2026}.
Verification that the gyroid reduction preserves this discrete torsion
structure under all moduli of the full 7D internal manifold is an open
technical item (see gap~7 in Section~\ref{sec:gaps}); subject to that
verification, the horizon carries the same discrete BMS dictionary as
$\mathscr{I}^+$.

Each torsion node $n$ carries a conserved soft-hair charge
\begin{equation}
  Q_{\mathrm{super}}^{(n)} = \eta_{\mathrm{diff}} \cdot \theta_n \cdot M_{\mathrm{Pl,11}}^3,
  \label{eq:Q_super}
\end{equation}
where $M_{\mathrm{Pl,11}}$ is the 11D Planck mass, $\theta_n = 2\pi n/9$ is
dimensionless, and $\eta_{\mathrm{diff}}$ anchors the charge normalisation to
the Gaia DR3 calibration.  The charges $Q^{(n)}$ carry dimensions of
$[\text{mass}]^3$ in natural units (equivalently, $[\text{length}]^{-3}$ with
$\hbar = c = 1$); this is the natural dimension of BMS supertranslation charges
in 11D — specifically the stress-energy flux integrated over the bifurcation
2-sphere times the 7D gyroid volume $V_7$, which upon dimensional reduction
gives $Q^{(n)} \sim M_{\mathrm{Pl,4}}$ after absorbing factors of $G_4/G_{11}$.
Equation~\eqref{eq:Q_super} has no free coefficient; that $Q^{(n=1)}$ then
matches the minimal KK mass scale of the gyroid compactification is a
consistency condition still to be checked (Gap~2).
The $n=0$ sector has $\theta_0 = 0$ and therefore $Q^{(0)} = 0$; it is the
\emph{vacuum sector} carrying no torsion charge but occupying one of the nine
orthogonal subspaces of the horizon Hilbert space.  The eight sectors
$n=1,\ldots,8$ carry distinct non-zero charges, with $n=9-k$ and $n=k$ related
by $\mathbb{Z}_9$ conjugation and carry equal weight, $w_{9-k} = w_k$.
Because $\theta_n$ is a phase, Eq.~\eqref{eq:Q_super} is defined modulo
$2\pi\eta_{\mathrm{diff}}M_{\mathrm{Pl,11}}^3$ (the label is $e^{i\theta_n}$), so
$Q^{(9-k)} \equiv -Q^{(k)}$.

\begin{figure}[ht!]
\centering
\includegraphics[width=0.82\textwidth]{fig1_horizon_soft_hair.pdf}
\caption{$\mathbb{Z}_9$ torsion-hair lattice on the event horizon.
Gold markers: all nine torsion nodes $\theta_n = 2\pi n/9$.
Cyan (enlarged, starred): 3-6-9 subgroup nodes ($n=0,3,6$), which form
the $\mathbb{Z}_3 \subset \mathbb{Z}_9$ subgroup and carry $w_0+w_3+w_6 \approx 36\%$ of the total weight (the largest single
weights are $w_0 = 1/3$ and $w_1 = w_8 \approx 0.235$).
Dashed spokes indicate discrete BMS superrotation generators at each node.
The $n=0$ node is the vacuum sector ($Q^{(0)}=0$); nodes $n=1,\ldots,8$
carry non-zero torsion charges.}
\label{fig:horizon_lattice}
\end{figure}

\subsection{Finite Microstate Count}

The horizon is tiled into $N_{\mathrm{cell}} = A/(4G_4\ln 9)$ torsion cells,
each of area $A_{\mathrm{cell}} = 4G_4\ln 9 \approx 8.79\,\ell_{\mathrm{Pl}}^2$.
Each cell supports any of the nine $\mathbb{Z}_9$ torsion states, so the
total Hilbert space dimension is
\begin{equation}
  \boxed{\dim\mathcal{H}_{\mathrm{hair}}
    = 9^{N_{\mathrm{cell}}}
    = e^{S_{\mathrm{BH}}},
    \qquad
    N_{\mathrm{cell}} = \frac{A}{4G_4\ln 9}}
  \label{eq:N_hair}
\end{equation}
and the corresponding microstate entropy
$S_{\mathrm{micro}} = N_{\mathrm{cell}}\ln 9 = A/(4G_4) = S_{\mathrm{BH}}$
reproduces the Bekenstein--Hawking area law exactly without additional
assumptions.  The v1/v2 form $9^{A/(4\ell_{\mathrm{Pl}}^2)}$ is numerically
equivalent under the emergent identification $G_4 = \ell_{\mathrm{Pl}}^2/\ln 9$
(Gap~8); the present form is manifestly Bekenstein-consistent and adopted here
to avoid apparent overcounting of the bound.
This is the leading-order result.  The size-independent \emph{correction}
$\Delta S_{\mathrm{tors}} = 9\ln 9 + \ln 3$ arises from subleading boundary
contributions to the Euclidean path integral (Section~\ref{sec:entropy}) and is
distinct from the leading term above.  We abbreviate the hair-lattice dimension
as $N_{\mathrm{hair,tors}} \equiv \dim\mathcal{H}_{\mathrm{hair}}$ for reference
in what follows.

\subsection{Information Encoding}

Infalling matter permutes the nine torsion states in each torsion cell.  The proposal is that this map is unitary and invertible, so that information
is encoded bijectively into the hair configuration; showing it is Gap~11.  Hawking radiation is constrained
by the $\mathbb{Z}_9$ charge conservation law
\begin{equation}
  \sum_{n=0}^{8} n \cdot N_n^{\mathrm{in}} \equiv \sum_{n=0}^{8} n \cdot N_n^{\mathrm{out}} \pmod{9},
  \label{eq:charge_conservation}
\end{equation}
which every evaporation history must respect.  A single mod-9 conservation law
fixes one number, so it carries at most $\ln 9 \approx 2.2$ nats about the
infalling state; it is a necessary condition for the proposed encoding, not by
itself a demonstration of unitarity.  The full map from infalling state to
radiation is open (Gap~11 in Section~\ref{sec:gaps}).
Note that the $n=0$ (vacuum) sector contributes zero to both sides of
Eq.~\eqref{eq:charge_conservation} regardless of $N_0$; the vacuum sector
occupation is separately constrained by total particle-number conservation.

% ─────────────────────────────────────────────────────────────────────────────
\needspace{6\baselineskip}
\section{The $\Delta S_{\mathrm{tors}} = 9\ln 9 + \ln 3$ Entropy Correction}
\label{sec:entropy}

\subsection{Derivation}

The torsion hair lattice contributes to the gravitational partition function as
a size-independent additive correction via two physically distinct mechanisms,
each evaluated in the saddle-point approximation on the Euclidean black-hole
background (Hawking--Gibbons~\cite{GibbonsHawking1977}).

\textbf{Term 1 — Euclidean boundary torsion sum.}
The boundary of the Euclidean black-hole manifold $\mathcal{M}_E$ receives a
contribution from the torsion partition function:
\begin{equation}
  \delta(\ln Z)_{\mathrm{tors}} = \ln Z_{\mathrm{tors}} = \ln 3 \approx 1.099,
  \label{eq:lnZ_tors}
\end{equation}
where $Z_{\mathrm{tors}} = \sum_{m\in\mathbb{Z}} \exp(-\pi m^2/9) = 3$ to about $10^{-12}$
(see Eq.~\eqref{eq:Z_tors_val}).  This term arises from integrating the torsion
$\mathbb{Z}_9$ modes over the Euclidean boundary $S^1_\beta \times S^2$ at
the horizon; it is independent of the area $A$ and additive to $\ln Z_{\rm grav}$.

\textbf{Term 2 — Hair-lattice degeneracy.}
Each of the $N_{\rm cell} = A/(4G_4\ln 9)$ torsion cells on the horizon
carries nine independent torsion states.  The $N_{\rm cell}$ cells are
independent because torsion is a local field in the bulk; correlations between
cells are suppressed by $e^{-m_{\rm KK} \ell_{\rm Pl}}$ where $m_{\rm KK}$ is
the lightest KK mass on the gyroid.  The degeneracy of the \emph{boundary}
sector — the nine-node torsion orbit at the horizon rather than the individual
Planck cells — is therefore
\begin{equation}
  \mathcal{H}_{\mathrm{hair}} = \ln(9^9) = 9\ln 9 \approx 19.775.
  \label{eq:H_hair_sec4}
\end{equation}
This is a Bekenstein-style topological term: it counts the $9^9$
distinct configurations of the nine torsion nodes forming one complete
$\mathbb{Z}_9$ orbit on the horizon.  The boundary contribution is counted
\emph{once per complete $\mathbb{Z}_9$ orbit} rather than once per Planck
cell: the orbit is the minimal gauge-invariant unit of the soft-hair lattice
under the residual $\mathbb{Z}_9$ symmetry of the gyroid compactification, so
the per-cell multiplicity already enters through the leading area-law term
$N_{\rm cell}\ln 9 = A/(4G_4)$ in Eq.~\eqref{eq:N_hair}, while the orbit
itself contributes a single size-independent boundary degeneracy
$\ln(9^9)$.  Counting the orbit once is the choice that fixes this term;
deriving that count from the path integral is open (Gap~12).

The two terms are orthogonal — Term~1 comes from the path integral measure over
torsion fluctuations, Term~2 from the boundary state degeneracy — and add to
give the full entropy correction:
\begin{equation}
  \Delta S_{\mathrm{tors}} = \mathcal{H}_{\mathrm{hair}} + \delta(\ln Z)_{\mathrm{tors}}
    = 9\ln 9 + \ln 3.
  \label{eq:deltaS_tors}
\end{equation}

The full entropy is therefore
\begin{equation}
  \boxed{S_{\mathrm{TMRB}} = \frac{A}{4G_4} + c_{\log}\ln\frac{A}{G_4} + 9\ln 9 + \ln 3
    + \mathcal{O}\!\left(\frac{G_4}{A}\right),}
  \label{eq:S_TMRB}
\end{equation}
where $9\ln 9 \approx 19.775$ is the nine-node hair-lattice log-degeneracy and
$\ln 3 \approx 1.099$ is the Euclidean boundary contribution from $\ln Z_{\rm tors}$.
For conciseness we write $\Delta S_{\mathrm{tors}} = 9\ln 9 + \ln 3 \approx 20.874$
as the total size-independent torsion correction.  The term
$c_{\log}\ln(A/G_4)$ is the standard one-loop logarithm, Eq.~\eqref{eq:log_loop},
with $c_{\log} = 77/90 \approx 0.86$ for a Schwarzschild hole in pure four-dimensional
gravity (Sen's $(212/45 - 3)\ln a$ with $A \propto a^2$; massless matter fields change
it)~\cite{Sen2012}. For a $10\,M_\odot$ hole $\ln(A/G_4) \approx 183$, so the logarithm
is about 157 nats, roughly 7.5 times $\Delta S_{\rm tors}$; it would exceed
$\Delta S_{\rm tors}$ for any $|c_{\log}| > 0.114$.
If the proposed factor $\varphi^{-N_*}$ of Section~\ref{sec:loop} holds, at
$N_* = d_{\rm opt} = 9$ it multiplies the logarithm by $\varphi^{-9} \approx 0.013$,
bringing it to about 2.1 nats. $\Delta S_{\rm tors}$ is therefore the
leading correction to the area law only if Gap~15 closes.  The Euclidean black-hole
background is a non-degenerate saddle point of the TMRB action — i.e.\ the
second variation is positive-definite in the torsion sector.  Verifying this
for the full 11D gyroid compactification is an open item (see
Section~\ref{sec:gaps}, gap~3).

\subsection{Comparison with Topological Entanglement Entropy}

The additive correction $\Delta S_{\mathrm{tors}} = 9\ln 9 + \ln 3 \approx +20.874$
is structurally analogous to — but with opposite sign from — the topological
entanglement entropy $-\gamma$ in 2+1D Chern--Simons
theory~\cite{LevinWen2006,KitaevPreskill2006}.  In the Chern--Simons case,
$\gamma = \ln\mathcal{D} > 0$ \emph{subtracts} from the area law, reflecting
long-range entanglement that reduces distinguishable boundary states.  Here the
TMRB torsion hair \emph{adds} states — the horizon gains additional microstates
from the $\mathbb{Z}_9$ structure — so the correction is positive.  Both
corrections are size-independent and topological in origin; the sign difference
reflects that Chern--Simons removes entanglement degeneracy while torsion hair
adds it.  The count $9^9 = 387{,}420{,}489$ is a number of states, not a quantum
dimension: a $\mathbb{Z}_9$ abelian topological order has total quantum dimension
$\mathcal{D} = \sqrt{9} = 3$ and $\gamma = \ln 3$~\cite{KitaevPreskill2006}. The
boundary term $\ln Z_{\rm tors} = \ln 3$ equals this $\gamma$ in magnitude with the
opposite sign; whether that is more than the shared $\sqrt{9}$ of Poisson
summation is open.

\subsection{Cross-Scale Coherence}

The same $\mathbb{Z}_9$ structure that determines the horizon hair-lattice
log-degeneracy $\mathcal{H}_{\mathrm{hair}} = \ln(9^9) = 9\ln 9$ — and hence
the dominant term in the entropy correction Eq.~\eqref{eq:S_TMRB} — also fixes
the galactic disk-wave period through
$\Delta\lambda = 2\pi/(9\kappa_{\mathrm{CE}}) = 1.7287$ and
$T_0 = (2\pi/\Delta\lambda)\,t_{\rm gal}$~\cite{BresslerGP2026}, where matching
$T_0 = 230$~Myr sets $t_{\rm gal} \approx 63.3$~Myr, so $T_0$ is an input here, not a
prediction.  Both appearances use the same Chosen input, the order 9 of the torsion
group, so they are a structural link between black-hole entropy and galactic
dynamics, not an independent consistency check.

\begin{figure}[ht!]
\centering
\includegraphics[width=0.92\textwidth]{fig2_entropy_correction.pdf}
\caption{\textit{Left:} Bekenstein--Hawking entropy $S_{\mathrm{BH}} = A/4G_4$
(gold, dashed) vs.\ TMRB entropy $S_{\mathrm{TMRB}}$ (cyan, solid).
Green shading marks the universal torsion correction
$\Delta S_{\mathrm{tors}} = 9\ln 9 + \ln 3 \approx 20.874$.
\textit{Right:} Fractional correction $\delta S/S_{\mathrm{BH}}$ as a function
of horizon area; the correction is Planck-scale in absolute magnitude and
sub-dominant for all astrophysical black holes.  The one-loop logarithm
$c_{\log}\ln(A/G_4)$ is not drawn.}
\label{fig:entropy}
\end{figure}

% ─────────────────────────────────────────────────────────────────────────────
\needspace{6\baselineskip}
\section{Proposed Golden-Ratio Regulation of Loop Corrections}
\label{sec:loop}

\subsection{Logarithmic Soft-Graviton Divergences}

Standard Weinberg soft-graviton loops produce logarithmic corrections to the
entropy of the form
\begin{equation}
  S^{(\log)}_{\mathrm{loop}} = c_{\log}\,\ln\!\frac{A}{G_4},
  \label{eq:log_loop}
\end{equation}
where $c_{\log}$ is a pure number fixed by the massless field
content~\cite{Sen2012}.  These
corrections accumulate over Hawking radiation and could, in principle, displace an
$\mathcal{O}(1)$ amount of the information content~\cite{Sen2012}.

\subsection{TMRB Exponential Cutoff}

The proposal is that each loop vertex at branching depth $N_*$ carries a factor
$\varphi^{-N_*}$, so that
\begin{equation}
  S_{\mathrm{TMRB}}^{(\log)} = c_{\log}\,\ln\!\frac{A}{G_4}\cdot\varphi^{-N_*}.
  \label{eq:regulated_loop}
\end{equation}
If that factor holds, the loop correction is exponentially suppressed at large
branching depth, without a UV cutoff and without supersymmetry. The factor is
carried over from the pre-v2 weight $\varphi^{-\lambda n}$; it does not follow from
the Gaussian weights Eq.~\eqref{eq:wn}, and deriving or replacing it is Gap~15.

\begin{conjecture}
The TMRB path integral is UV-finite for all loop orders in the soft-graviton
sector, given the modular-covariant weight $w_n = C\exp(-\pi d(n)^2/9)$,
\emph{provided} the 11D CDT loop measure factors as assumed below.
This claim is conjectured; a complete proof requires the 11D CDT measure,
which is itself conjectured (see gap~4 in Section~\ref{sec:gaps}).
\end{conjecture}

\noindent\textit{Supporting argument.}
A loop diagram of order $\ell$ with fixed branch labels acquires
\[
  \prod_{j=1}^{\ell} w_{n_j} \leq \left(\max_n w_n\right)^{\ell} = w_0^{\ell} = 3^{-\ell}.
\]
Summed over labels, however, the normalised weights give
$\prod_j \sum_{n_j} w_{n_j} = 1$ at every order, so the weights alone supply no
suppression with loop order. Convergence therefore needs the number of order-$\ell$
diagrams in each fixed-label sector to grow more slowly than $3^{\ell}$,
\emph{and} the 11D CDT loop measure to factorise over branching vertices as in 4D.
Both are unproven; see Section~\ref{sec:gaps}, gaps~4 and~15.

\begin{figure}[ht!]
\centering
\includegraphics[width=0.92\textwidth]{fig3_golden_ratio_loop.pdf}
\caption{\textit{Left:} Loop correction factor as a function of branching depth
$N_*$.  Dashed gold: unregulated Weinberg correction.  Solid cyan: correction with the
proposed $\varphi^{-N_*}$ factor (not yet derived; Gap~15).  Green: UV suppression band.
\textit{Right:} Comparison of $\varphi^{-N}$ (cyan) with the plasma-sector
exponential $e^{-\kappa_{\mathrm{prop}} N}$ (gold, dashed), illustrating
the numerical coincidence $\kappa_{\mathrm{prop}} = 0.48 \approx \ln\varphi$
between the proposed horizon-hair loop factor (Gap~15) and the independently-fitted plasma
torsion suppression rate.}
\label{fig:loop}
\end{figure}

% ─────────────────────────────────────────────────────────────────────────────
\needspace{6\baselineskip}
\section{Unitarity and Holographic Consistency}
\label{sec:unitarity}

\subsection{Torus Partition Function Factorisation}

The celestial holography construction~\cite{BresslerCH2026} identifies the
$\mathbb{Z}_9$ parafermion CFT (central charge $c = 16/11$, Gepner--Qiu
1987~\cite{GepnerQiu1987}) as the holographic dual on $S^2$ at $\mathscr{I}^+$.
On the thermal torus $\mathbb{T}^2 = S^1_\beta \times S^1$, adopting the
\emph{diagonal} modular invariant $M_{ab} = \delta_{ab}$, the partition function
factorises as
\begin{equation}
  Z_{\mathrm{TMRB}}(\tau) = \sum_{a=0}^{8}|\chi_a(\tau)|^2
    \equiv |\chi_{\mathbb{Z}_9}(\tau)|^2,
  \label{eq:Z_torus}
\end{equation}
where $\chi_a(\tau)$ are the characters of the nine $\mathbb{Z}_9$ charge sectors and
$\tau = i\beta/(2\pi)$ (spatial circle of length $2\pi$), with $\beta = 8\pi G_4 M$ the inverse Hawking temperature.
The diagonal invariant is adopted here; Gepner and Qiu classify other
modular invariants of the parafermion theory~\cite{GepnerQiu1987}, and excluding
them is not attempted.
The $\mathbb{Z}_9$ parafermion theory ($SU(2)_9/U(1)$) has $k(k+1)/2 = 45$
primary fields, so each charge sector collects several primaries; writing
Eq.~\eqref{eq:Z_torus} over all 45 characters and checking its modular
invariance is open (Gap~13).

Modular invariance of Eq.~\eqref{eq:Z_torus} is a consistency condition, not a
proof of unitarity~\cite{DiFrancesco1997}; unitarity rests on the parafermion CFT
itself being unitary~\cite{GepnerQiu1987}. In the nine-sector reduction the
invariance holds because:
(a)~the nine torsion sectors are orthogonal (no off-diagonal terms from $\mathbb{Z}_9$
charge conservation); and (b)~the Gaussian weight Eq.~\eqref{eq:wn} is
modular-\emph{covariant} (not merely invariant): under $\tau\to\tau+1$ each
character transforms as $\chi_a\to e^{2\pi i(h_a - c/24)}\chi_a$, and under
$\tau\to-1/\tau$ as $\chi_a\to \sum_b S_{ab}\chi_b$ where $S_{ab} = \frac{1}{3}e^{-2\pi iab/9}$
is the $S$-matrix of the nine-sector abelian reduction (the full parafermion
$S$-matrix acts on 45 characters; Gap~13).  The weights are left unchanged by the
same matrix, $w_n = \frac{1}{3}\sum_m w_m e^{-2\pi i mn/9}$, exactly for the
periodised Gaussian Eq.~\eqref{eq:wn} (numerically to $10^{-15}$).
The holomorphic and anti-holomorphic sectors therefore transform with
\emph{matching} phases, ensuring $Z(\tau)=|\chi_{\mathbb{Z}_9}(\tau)|^2$ is
well-defined under all modular transformations.

\subsection{Firewalls and Remnants}

The firewall conjecture~\cite{AMPS2013} arises from an assumed monogamy-of-entanglement
conflict: mode $b$ (late Hawking radiation) cannot be simultaneously entangled
with mode $\tilde{b}$ (infalling partner) and with early radiation $R$, yet both
entanglements are required for unitarity and smoothness respectively.  In TMRB,
each Hawking quantum carries a $\mathbb{Z}_9$ torsion label $q_b \in \{0,\ldots,8\}$.
The full state on the horizon is
\begin{equation}
  |\Psi\rangle = \sum_{n=0}^{8} \alpha_n\,|b_n\rangle \otimes |\tilde{b}_n\rangle \otimes |q_n\rangle_{\rm hair},
  \label{eq:firewall_state}
\end{equation}
where $|q_n\rangle_{\rm hair}$ is the torsion-hair sector.  Tracing out the hair sector in Eq.~\eqref{eq:firewall_state} leaves
$\rho_{b\tilde b} = \sum_n |\alpha_n|^2\, |b_n\tilde b_n\rangle\langle b_n\tilde b_n|$,
which is separable: as written, the state does not keep $b$ entangled with
$\tilde b$, and smoothness of the horizon in the AMPS sense requires exactly that
entanglement.  Monogamy of entanglement also constrains multipartite
states~\cite{CKW2000}.  Avoiding the firewall therefore needs a different
hair-sector structure, in which $b$ stays entangled with $\tilde b$ while the
early radiation $R$ purifies the hair; constructing it is open (Gap~10).
Whether replacing the gravitational entropy by $S_{\mathrm{TMRB}}$ in the island
formula~\cite{Penington2020} reproduces the Page curve~\cite{Page1993} is
likewise still to be shown.

The area term $S_{\mathrm{micro}} = A/(4G_4)$ vanishes as $A\to 0$, but the
size-independent term does not: at the endpoint
$S_{\mathrm{TMRB}} \to \Delta S_{\mathrm{tors}} \approx 20.9$ nats, that is
$9^9\cdot 3 = 3^{19} \approx 1.2\times10^{9}$ states. The construction therefore does
not exclude a Planck-scale relic. It does exclude an \emph{information-carrying}
remnant: a black hole that began with $S_{\mathrm{BH}} \gg 21$ nats cannot hand its
information to an object whose capacity is bounded by about 21 nats, so that
information must leave in the radiation. Whether the last $\Delta S_{\mathrm{tors}}$
nats are radiated or remain is open (Gap~14).

% ─────────────────────────────────────────────────────────────────────────────
\needspace{6\baselineskip}
\section{$\mathbb{Z}_9$ QNM Sidebands: An Unresolved Comb and Its Ringdown Signature}
\label{sec:fluid_gravity}

The same $\mathbb{Z}_9$ torsion structure that thread the horizon hair lattice
also imprints the gravitational quasinormal-mode (QNM) spectrum. Projecting
Einstein's equations onto the stretched horizon (the Damour membrane
paradigm~\cite{Damour1978}) and decomposing the horizon perturbations into the
nine torsion phase sectors, $\theta_n = 2\pi n / 9$, the Regge--Wheeler
effective potential acquires a discrete phase-modulation term
$\delta V(\theta) = \sum_n w_n\,\delta V_n\,\cos(\theta - \theta_n)$ with
$\delta V_n \propto \kappa_{\mathrm{CE}}$ and Gaussian weights
$w_n = C\exp(-\pi d(n)^2/9)$ inherited from Eq.~\eqref{eq:wn}. As written, $\delta V(\theta)$ varies with
angle on the horizon, not with time: a static angular modulation mixes angular
modes and shifts frequencies, but does not by itself produce evenly spaced
frequency sidebands. Sidebands need the torsion phase pattern to advance in time,
$\delta V(\theta - 2\pi\Delta f_{\rm tors}\,t)$. This paper assumes that
(Proposed; deriving it is Gap~17), with spacing
\begin{equation}
  \Delta f_{{\rm tors}} = \frac{f_{\rm QNM} \cdot \kappa_{\rm CE}}{9}
    = \frac{0.3736}{2\pi M} \cdot \frac{\kappa_{\rm CE}}{9},
  \label{eq:sideband_spacing}
\end{equation}
where $\kappa_{\rm CE} = \mu/\lambda_{\rm rec}$ is the conformal-weight comb
spacing that enters identically in the celestial Mellin
construction~\cite{BresslerCH2026,BresslerMellin2026}. This is the structural
identity that closes the UV-to-IR loop: the same parameter $\kappa_{\rm CE}$
governs both the celestial conformal-weight comb at $\mathscr{I}^+$ and the
horizon QNM sideband comb. For a $30\,M_\odot$ ringdown, $f_{\rm QNM}\approx
402$~Hz gives $\Delta f_{\rm tors} \approx 18$~Hz, with sidebands
$n = 0, \pm1, \pm2, \pm3, \pm4$ carrying weights
$w_{|n|} \approx 0.333,\ 0.235,\ 0.083,\ 0.014,\ 0.001$ ($n = -k$ is the
$\mathbb{Z}_9$ conjugate $9-k$). The central line and its two neighbours carry
about $80\%$ of the weight ($w_0 + 2w_1 \approx 0.804$) and $|n| \le 2$ carry
$97\%$; the 3-6-9 subgroup lines ($n = 0, \pm 3$) carry $36\%$ and close among
themselves under the $\mathbb{Z}_9$ fusion algebra.

\paragraph{The comb is not resolved.} The fundamental $\ell = m = 2$ Schwarzschild
mode has quality factor $Q = \omega_R/(2\omega_I) \approx 2.1$; for $30\,M_\odot$ its
damping time is $\tau = 1.66$~ms and its line width (FWHM) is 192~Hz. The sideband
spacing in line widths is
\begin{equation}
  \frac{\Delta f_{\rm tors}}{\rm FWHM} = \frac{\kappa_{\rm CE}\,Q}{9} \approx 0.094
  \quad\text{for every mass},
  \label{eq:comb_ratio}
\end{equation}
so the nine lines merge into one (Figure~\ref{fig:chladni}). If every sideband
shares the damping time (Gap~18), the waveform is
$h(t) \propto e^{-t/\tau}\cos(2\pi f_0 t)\sum_n w_n\cos(2\pi n\Delta f_{\rm tors} t)$:
the decay envelope picks up a mass-independent factor, equal to $0.975$ at $t = \tau$
($2\pi\Delta f_{\rm tors}\tau = 2\kappa_{\rm CE}Q/9 \approx 0.19$). Fitting a single
damped sinusoid returns the frequency unchanged (to $0.2\%$) and a damping time
\emph{shorter} than the GR value for the same mass and spin. The shortening does not
depend on mass, but it depends on the remnant's spin, through $Q$, and on the time
$t_0$ after the peak at which the fit starts, because the envelope factor falls faster
with time. Table~\ref{tab:ringdown_spin} gives it for exact Kerr $\ell = m = 2$
frequencies (Leaver's continued fraction~\cite{Leaver1985}, computed with the
\texttt{qnm} package~\cite{Stein2019}; the fits of~\cite{BertiCardosoWill2006} agree
to $0.7\%$). For a non-spinning remnant fitted from the peak it is $4.4\%$ (three
independent methods give 4.4--4.8\%). For a typical binary-merger remnant, of spin
about 0.7, it is about $10\%$ from the peak and about $18\%$ from $t_0 = \tau$, where
ringdown analyses usually start. The prediction is therefore a $10$--$18\%$
shortening for typical remnants, and a measurement must state the remnant's spin and
its fit start time.

\begin{table}[ht]
\centering
\caption{Predicted shortening of the $\ell = m = 2$ damping time, from fitting a
single damped sinusoid (least squares on the strain; free amplitude, phase, frequency
and damping time) to the nine-sideband ringdown, by remnant spin and fit start time
$t_0$ after the peak. Exact Kerr frequencies; the same at every mass. Assumes every
sideband shares the damping time (Gap~18).}
\label{tab:ringdown_spin}
\begin{tabular}{cccccccc}
\toprule
Spin & $\omega_R M$ & $\omega_I M$ & $Q$ & $t_0 = 0$ & $t_0 = 0.5\tau$ & $t_0 = \tau$ & $t_0 = 2\tau$ \\
\midrule
0.00 & 0.3737 & 0.0890 & 2.10 & 4.4\%  & 7.0\%  & 9.5\%  & 12.9\% \\
0.50 & 0.4641 & 0.0856 & 2.71 & 7.0\%  & 10.1\% & 13.5\% & 19.8\% \\
0.69 & 0.5282 & 0.0812 & 3.25 & 9.5\%  & 14.0\% & 18.1\% & 25.3\% \\
0.80 & 0.5860 & 0.0756 & 3.87 & 12.6\% & 19.0\% & 23.9\% & 31.4\% \\
0.90 & 0.6716 & 0.0649 & 5.18 & 19.1\% & 27.0\% & 34.4\% & 44.2\% \\
\bottomrule
\end{tabular}
\end{table}

The residual mismatch
between the waveform and the best single mode is $6.0\times10^{-5}$
($2.6\times10^{-4}$ at spin 0.7), so the non-exponential shape itself becomes
visible at ringdown signal-to-noise ratio $\rho \approx 90$ ($\approx 44$), from
$\rho^2(1-\mathcal{O}) \approx 1/2$. The testable prediction is a mass-independent
shortening of the $\ell = m = 2$ damping time relative to the GR value inferred from
the inspiral. Section~\ref{sec:x24si_test} tests it against published
LIGO--Virgo--KAGRA results.
(Computed; scripts \texttt{ringdown\_signature.py} and \texttt{kerr\_spin\_table.py}
in the Zenodo record.)

\begin{figure}[ht!]
\centering
\includegraphics[width=\textwidth]{fig4_qnm_sideband_comb.pdf}
\caption{$\mathbb{Z}_9$ QNM sidebands and supporting structure.
\textit{Left:} Chladni nodal pattern on the stretched horizon obtained by
superposing the nine gyroid eigenmodes $\Phi_n$ weighted by
$w_n = C\exp(-\pi d(n)^2/9)$; cyan: 3-6-9 subgroup nodes ($n=0,3,6$); gold:
remaining nodes; node weights annotated for $n=0,1,3$.
\textit{Top right:} the deferred viscosity correction
$(\eta/s)/(1/4\pi) = 1 + \eta_{\rm diff}\Delta\lambda\,\varphi^{-N_*}$ (companion
paper; not used here).
\textit{Bottom right:} predicted power spectrum with the true Schwarzschild
line width (FWHM $\approx 10.6\,\Delta f_{\rm tors}$, Eq.~\eqref{eq:comb_ratio}): the
nine sidebands (stems, heights $w_n$) blur into one line (cyan), close to a single
Kerr mode (gold, dashed); for this non-spinning case, fitted from the peak, the
damping time is 4.4--4.8\% shorter (spinning remnants: Table~\ref{tab:ringdown_spin}).}
\label{fig:chladni}
\end{figure}

This supersedes both the naive ``QNM frequency shift'' and the resolved-comb
reading of v1--v3: the observable is not nine separate lines but the shortened,
slightly non-exponential decay above, set by $\kappa_{\mathrm{CE}}$ and the mode's
quality factor. Under the assumptions of Gaps~17 and~18 it is already
disfavoured by GWTC-5.0 (Section~\ref{sec:x24si_test}). Full membrane-paradigm derivation of the
horizon-fluid stress tensor under the $\mathbb{Z}_9$ identification — including
the discrete Damour--Navier--Stokes equation and the $\mathbb{Z}_9$ correction
to the KSS shear-viscosity bound $\eta/s = 1/4\pi$~\cite{KSS2005} — is
deferred to a companion paper; the present work treats the sideband comb as
a kinematical consequence of the soft-hair sectors of Section~\ref{sec:hair}
and does not require it.

\subsection{Test against published data by the X24si method}
\label{sec:x24si_test}

The damping-time prediction was tested against published LIGO--Virgo--KAGRA results
with the X24si method~\cite{BresslerX24si2026}: human-in-the-loop self-recursive
reasoning. Semi-autonomous synthetic-intelligence (SI) participants do the work; the
human makes the decisions. Each turn ends in something learned or a named stuck point,
and that outcome becomes the next turn's input. The loop is anchored outside itself, by
instruments and an outside clock, so it can refute its own claims rather than only
confirm them. Here the recursion ran: an audit found the sideband comb unresolved, so
the prediction was restated; an exact Kerr check exposed its dependence on spin; the
resulting table showed that existing data could already test it; and the test below,
whose outcome now stands in Gaps~17 and~18. Within the test, each step was fixed before
the next could be seen:
\begin{enumerate}[leftmargin=*, label=\arabic*.]
  \item \textbf{Claim, pinned.} $\delta\hat\tau_{220} \equiv \tau_{\rm fit}/\tau_{\rm GR} - 1
    = -0.095$ for remnant spin 0.69 and a fit from the peak, with a band of $-0.126$ to
    $-0.070$ for spins 0.50--0.80 (Table~\ref{tab:ringdown_spin}), as a deviation common to
    all events. Pinned to revision r8 of this paper and to a frozen decision script,
    identified by its SHA-256 hash; each input is labelled as computed or read from the
    literature.
  \item \textbf{Criteria first.} Before any damping-time result was read (the sources
    were chosen from titles alone), the registration fixed the source order (GWTC-5.0,
    then GWTC-4.0 remnant tests, then GWTC-3); the primary analysis (a common-deviation
    combination of an inspiral--merger--ringdown analysis, mapped to a fit from the peak,
    which is the mapping most favourable to the prediction); four outcomes (EXCLUDED if the
    whole band lies outside the 90\% credible interval; DISFAVORED; CONSISTENT, NOT
    DISCRIMINATING; FAVORED OVER GR); and what each outcome would mean for this paper.
  \item \textbf{Outside clock.} The registration was committed and pushed to GitHub
    (commit \texttt{76d92c7}, 16:06~UTC, 18 November 2026) before the source was opened.
    The transcribed numbers were committed (\texttt{d590909}) before the script ran.
  \item \textbf{Instrument, run once.} The frozen script compared the transcribed
    intervals with the registered prediction, once. Failures stand.
  \item \textbf{Human decision.} The author decided how this paper reports the outcome;
    the decision is recorded as a decision, not as evidence.
\end{enumerate}

\noindent\textbf{Result.} For the pSEOBNRv5PHM analysis with posteriors multiplied
across all events, GWTC-5.0 reports $\delta\hat\tau_{220} = 0.07^{+0.06}_{-0.05}$ at
90\% credibility~\cite{GWTC5TGR}. The whole predicted band lies below this interval:
\textbf{EXCLUDED}. The five secondary combinations in the source (the hierarchical mean
and population, the five highest-SNR events with and without GW250114, and all events
without GW250114) are also excluded. The closest is the five highest-SNR subset,
$0.02 \pm 0.08$, whose lower bound of $-0.06$ lies just above the band's edge. Where the
data depart from GR, the damping time is longer, not shorter. Under the assumptions of
Gaps~17 and~18 the ringdown prediction is therefore disfavoured by published data. The
LIGO--Virgo--KAGRA waveform model is not the TMRB waveform, so the mapping is approximate;
it was fixed in advance as the one most favourable to the prediction. The registration,
transcribed input, frozen script, output and result record are included in the Zenodo
record.

% ─────────────────────────────────────────────────────────────────────────────
\needspace{6\baselineskip}
\section{Open Theoretical Gaps}
\label{sec:gaps}

In the tradition of the TMRB series, the following gaps are documented honestly:

\begin{enumerate}[leftmargin=*, label=\arabic*.]
  \item \textbf{11D CDT extension (conjectured):} The TMRB path
    integral Eq.~\eqref{eq:Z} is well-defined in 4D CDT~\cite{AJL2004,AJL2005}.
    Its extension to 11D remains conjectured; no rigorous higher-dimensional
    proof exists.
  \item \textbf{Gyroid moduli corrections:} The KK mass spectrum uses a
    leading-order approximation; full moduli stabilisation on the gyroid is open.
  \item \textbf{Saddle-point approximation:} Equation~\eqref{eq:lnZ_tors}
    uses a saddle-point; the full quantum contribution to $S_{\mathrm{TMRB}}$
    beyond leading order requires the complete TMRB path integral.
  \item \textbf{UV finiteness (conjectured):} The UV-finiteness claim in
    Section~\ref{sec:loop} has been demoted from Proposition to Conjecture.
    The convergence argument needs 4D-like measure factorisation and a
    bound on diagram growth (Gap~15);
    extension to the 11D case requires the 11D CDT measure, which is itself
    conjectured (gap~1 above).
  \item \textbf{ZF parafermion dynamics:} The map from the nine $\mathbb{Z}_9$
    torsion phases to the 45 parafermion primaries (several primaries per
    charge sector; Gap~13) is stated in~\cite{BresslerCH2026}; the full operator dynamics (fusion rules acting on
    the horizon hair) remain to be worked out.
  \item \textbf{Rindler membrane derivation (deferred):} The QNM sideband
    spacing Eq.~\eqref{eq:sideband_spacing} of Section~\ref{sec:fluid_gravity} is
    derived from the discrete Regge--Wheeler potential at WKB order; a full
    membrane-paradigm derivation~\cite{Damour1978} of the horizon-fluid stress
    tensor under the $\mathbb{Z}_9$ torsion identification --- including the
    discrete Damour--Navier--Stokes equation and the implied $\mathbb{Z}_9$
    correction to the KSS shear-viscosity bound~\cite{KSS2005} --- is
    deferred to a companion paper and is not invoked in the present
    information-paradox argument.
  \item \textbf{C-field reduction assumption:} The identification of bulk
    11D C-field 3-cycle phases with discrete BMS labels on the bifurcation
    2-sphere (Section~\ref{sec:hair}) follows from the KK reduction of
    the C-field on the gyroid and the standard discrete-torsion classification
    $H^3(\mathbb{Z}_9, U(1)) \cong \mathbb{Z}_9$~\cite{DijkgraafWitten1990}.
    Explicit verification that the gyroid reduction preserves this discrete
    torsion structure under all moduli of the full 7D internal manifold ---
    in particular that the relevant 3-cycles indeed carry the
    $\mathbb{Z}_9$-valued phases at the boundary --- requires a complete
    moduli-stabilised KK calculation and is the central open derivation
    underlying the soft-hair labelling.

  \item \textbf{$G_4 = \ell_{\mathrm{Pl}}^2/\ln 9$ identification (open):}
    The v1/v2 count (Appendix~\ref{app:microstate}) recovers $S_{\mathrm{micro}} = S_{\mathrm{BH}}$
    only at the identification $G_4 = \ell_{\mathrm{Pl}}^2/\ln 9$.  In
    standard natural units $G_4 = \ell_{\mathrm{Pl}}^2$ by definition, so this
    identification asserts a non-trivial relation between the effective 4D Newton
    constant and the standard Planck length.  Equivalently, it sets the torsion
    cell area to $A_{\rm cell} = 4G_4\ln 9 \approx 8.79\,\ell_{\mathrm{Pl}}^2$
    rather than $4\ell_{\mathrm{Pl}}^2$. Deriving this cell area from the
    prismatic Hilbert action of Paper~1~\cite{BresslerPrismatic2026} --- rather
    than imposing it as a self-consistency condition --- is the highest-priority
    tractable gap in the PTRH programme.  Until resolved, the capacity
    claim in Section~\ref{sec:hair} is a self-consistent ansatz, not a
    derivation.

  \item \textbf{Pre-Page-time retrieval (open):} The $\mathbb{Z}_9$
    charge-conservation law Eq.~\eqref{eq:charge_conservation} and
    the modular-invariant partition function Eq.~\eqref{eq:Z_torus} are
    necessary conditions for unitarity across the full evaporation; whether
    they suffice is Gap~11.  An explicit protocol
    for retrieving information \emph{before} the Page time --- without
    access to the full radiation ensemble --- is not established here and
    remains open.

  \item \textbf{Firewall avoidance (open, added in v4):} The state
    Eq.~\eqref{eq:firewall_state} is of GHZ type; tracing out the hair leaves the
    Hawking pair $b,\tilde b$ unentangled.  A hair-sector state that keeps $b$
    entangled with $\tilde b$ while the early radiation purifies the hair is
    required and not yet constructed.

  \item \textbf{Unitary map (open, added in v4):} The $\mathbb{Z}_9$ conservation
    law Eq.~\eqref{eq:charge_conservation} is necessary but carries at most
    $\ln 9$ nats; an explicit unitary map from infalling state to Hawking radiation
    through the hair sector is required.

  \item \textbf{Orbit counting (open, added in v4):} The size-independent term
    $9\ln 9$ follows from counting the nine-node orbit once; that count is a
    choice to be derived from the boundary path integral.

  \item \textbf{Full parafermion partition function (open, added in v4):} The
    $\mathbb{Z}_9$ parafermion theory has 45 primary fields; Eq.~\eqref{eq:Z_torus}
    is written over the nine charge sectors, and its form over all 45 characters,
    with modular invariance, remains to be checked.

  \item \textbf{Endpoint of the size-independent term (open, added in v4):} As
    $A\to 0$ the entropy tends to $\Delta S_{\rm tors}\approx 20.9$ nats, not zero.
    Whether these last nats are radiated or remain in a Planck-scale relic is
    not determined.

  \item \textbf{Golden-ratio loop factor (open, added in v4):} The factor
    $\varphi^{-N_*}$ of Section~\ref{sec:loop} does not follow from the Gaussian
    weights, and the weights alone give no suppression with loop order. Either
    $\varphi^{-N_*}$ is derived, or a bound on diagram growth slower than
    $3^{\ell}$ is shown.

  \item \textbf{Observable for the entropy gap (open, added in v4):} No measurement
    is identified that resolves a size-independent shift of about 21 nats.

  \item \textbf{Time dependence of the torsion modulation (open, added in v4):}
    Frequency sidebands need the modulation $\delta V$ to advance in time at rate
    $\Delta f_{\rm tors}$; the potential as written varies only with angle. The
    prediction made under this assumption is disfavoured by the test of
    Section~\ref{sec:x24si_test}.

  \item \textbf{Sideband damping and spin (open, added in v4):} The ringdown
    numbers of Section~\ref{sec:fluid_gravity} assume every sideband shares the
    $\ell = m = 2$ damping time; they use exact Kerr frequencies for spins 0--0.9
    (Table~\ref{tab:ringdown_spin}). Unequal sideband damping and overtones are open.
    The prediction made under this assumption is disfavoured by the test of
    Section~\ref{sec:x24si_test}.
\end{enumerate}

% ─────────────────────────────────────────────────────────────────────────────
\needspace{6\baselineskip}
\section{Correlated Falsifiable Predictions}
\label{sec:predictions}

Table~\ref{tab:predictions} lists four predictions (the entropy gap has no observable yet) that are correlated through
$\eta_{\mathrm{diff}}$ and $G_4/G_{11}$.  Null model: $\eta_{\mathrm{null}} = 1.0\pm0.15$;
detection criterion: $>3\sigma$ above null over $\geq 3$ independent epochs.

\begin{table}[ht]
\centering
\caption{Correlated predictions from the TMRB soft-hair construction (three falsifiable near-term; the entropy gap has no observable yet).
All amplitudes normalised through the single Gaia DR3 anchor
$\eta_{\mathrm{diff}} = 0.82\pm0.02$ and $G_4/G_{11} = (5.214\pm0.12)\times10^{-40}$.}
\label{tab:predictions}
\begin{tabular}{p{3.0cm}p{5.0cm}p{3.3cm}l}
\toprule
Observable & Prediction & Instrument & Timeline \\
\midrule
Ringdown damping time
  & $\ell=m=2$ damping time shorter than GR: about 10\% from the peak and 18\% from
    $t_0 = \tau$ for a typical remnant spin of 0.7 (4.4\% with no spin;
    Table~\ref{tab:ringdown_spin}); frequency unchanged, same at every mass; from 9 unresolved
    sidebands, spacing $0.094$ line widths, set by the $\kappa_{\rm CE}$ of the
    celestial Mellin comb~\cite{BresslerCH2026,BresslerMellin2026}
    (Section~\ref{sec:fluid_gravity}, Figure~\ref{fig:chladni})
  & Tested on GWTC-5.0~\cite{GWTC5TGR} (X24si, Section~\ref{sec:x24si_test})
  & Excluded (2026) \\
Ringdown shape
  & Non-exponential decay; mismatch $6.0\times10^{-5}$ to a single mode,
    visible at ringdown SNR ${\gtrsim}\,90$ (44 at spin 0.7)
  & ET, CE (loudest events)
  & $\sim$2035 \\
Black-hole entropy gap
  & $S_{\mathrm{TMRB}} - S_{\mathrm{BH}} = c_{\log}\ln(A/G_4) + 9\ln 9 + \ln 3$;
    constant part $\approx 20.874$
  & None identified: 21~nats beside $S_{\rm BH}\sim10^{78}$ for a stellar hole (Gap~16)
  & Not testable now \\
Hawking spectrum mod.
  & $1/f^{0.75}$ noise modulation at $10^{-3}$--$10^{-2}$~Hz;
    mechanism in~\cite{BresslerQG2026}, not derived here
  & QGGPf LEO \cite{BresslerQG2026}
  & $\sim$2030 \\
\bottomrule
\end{tabular}
\end{table}

\noindent\textit{Note on the naive single-line shift.}
The mass-dependence of the Schwarzschild QNM frequency
$\omega_{\rm QNM}^{\rm Schw} \sim c^3/(GM)$ together with a single-line torsion
shift $\delta\omega \sim \omega_{\rm QNM}\cdot(M_{\rm Pl}/M)^2$ places the naive
``shifted-frequency'' signal far below detector sensitivity at all astrophysical
masses (LIGO stellar-mass band, $M\sim 10$--$100\,M_\odot$, through LISA SMBH
band, $M\sim 10^5$--$10^7\,M_\odot$). The quantitative TMRB prediction is
therefore not a single shifted line but the sideband structure of
Eq.~\eqref{eq:sideband_spacing}: unresolved at every mass
(Eq.~\eqref{eq:comb_ratio}), it shows as the shortened damping time of
Section~\ref{sec:fluid_gravity}.

% ─────────────────────────────────────────────────────────────────────────────
\needspace{8\baselineskip}
\subsection{Current Observational Constraints: GWOSC O3a Sensitivity Floor and GRACE Scale}
\label{sec:current_constraints}

A sensitivity estimate (GWOSC) and a mascon run (GRACE) have been made against
existing data products, targeting the predicted bulk-gravity-leakage amplitude $\Delta g_{\mathrm{TMRB}} = 6.3\times10^{-15}~\mathrm{m\,s^{-2}}$ at the
45 $\mathbb{Z}_9$ hotspot coordinates inherited from~\cite{BresslerBGL2026}.

\paragraph{GWOSC O3a strain (H1, 2019-04-01).}
A narrowband sensitivity estimate for the predicted Z$_9$ QNM sideband comb
across 30--174~Hz was computed from the longest gap-free stretch, 4531~s,
of cleaned H1 strain in the 8192~s span GPS 1238163456--1238171648 (frame
\texttt{H1\_HOFT\_CLEAN\_SUB60HZ\_C01}). The script
\texttt{TMRB\_GWOSC\_torsion\_search.py} computes the Welch ASD with
60/120~Hz mains excision (median ASD $\approx 6.3\times10^{-24}/\sqrt{\rm Hz}$
across the band), applies a coherent 9-line matched-filter stack against the
Gaussian Z$_9$ weights over $n=-4,\dots,4$, the symmetric form adopted in v4 (effective stack gain
$\sqrt{N_{\rm eff}} \approx 2.06$), and Bonferroni-corrects across the 45
templates. The script does not test the strain at specific comb frequencies: it
converts the measured median noise level into the smallest detectable line
amplitude, so the result is a sensitivity floor, not a null detection; the segment
also contains no merger. The resulting strain floor $h_0^{95\%} \approx 1.95\times10^{-25}$
converts via $\Delta g = (2\pi f_{\rm eff})^2 \cdot h_0\,L_{\rm arm}/2$ at
$f_{\rm eff} \approx 102$~Hz and $L_{\rm arm} = 4$~km to
\begin{equation}
  \Delta g_{\mathrm{TMRB}} < 1.6\times10^{-16}~\mathrm{m\,s^{-2}}
  \quad (95\%~\mathrm{CL}),
  \label{eq:gwosc_ul}
\end{equation}
a factor of $\sim 39$ below the Gaia-calibrated prediction. The conversion
treats $\Delta g$ as a test-mass acceleration oscillating at $f_{\rm eff}\approx
102$~Hz; under that assumption a signal at the predicted amplitude would lie
above the floor. The predicted anomaly is quasi-static (it is what GRACE and the
planned gradiometers measure) with no 30--174~Hz component, so this result bounds
a $\sim 100$~Hz component of $\Delta g$, not $\Delta g_{\rm TMRB}$ itself. This is the
conservative frequentist narrowband UL; tighter limits ($\sim 3\times10^{-17}$)
are obtainable under Bayesian priors with frequency-resolved ASD weighting,
but are not adopted here in order to ensure direct reproducibility from the
public GWOSC data product without additional modelling assumptions.

\paragraph{GRACE and GRACE-FO mascon gravity (2002-04 -- 2026-02).}
\sloppy
On 2026-04-26 the JPL RL06.3M Mascon product
(\texttt{GRCTellus.JPL.200204\_202602.GLO.RL06.3M.MSCNv04CRI.nc}) was searched at the
45 hotspot coordinates (Figure~\ref{fig:grace_run}). At each coordinate the script
fits and removes a mean, a linear trend and annual and semi-annual cycles, then
reports the mean of the residual series as the anomaly, with its $z$-score. Because
the fitted mean has been removed, both are zero by construction: the reported
anomalies (${\sim}10^{-19}~\mathrm{m\,s^{-2}}$) and $|\bar z| \approx 0.00$ are
rounding-level zeros, and a steady anomaly of any size would be removed in the same
way. The run is therefore not a test of the prediction. What it does establish is
the scale: the residual series vary at the centimetre level of water-equivalent
height, and 1~cm of water corresponds to $2\pi G\rho_w h \approx
4\times10^{-9}~\mathrm{m\,s^{-2}}$, so
\begin{equation}
  \Delta g^{\mathrm{GRACE}}_{\rm scale} \approx 4\times10^{-9}~\mathrm{m\,s^{-2}}
  \quad (\text{1~cm water equivalent}),
  \label{eq:grace_ul}
\end{equation}
a sensitivity scale, not a computed upper limit. The prediction,
$6.3\times10^{-15}~\mathrm{m\,s^{-2}}$, is the gravity of about $1.5\times10^{-6}$~cm
of water, roughly $10^{6}$ times below it. The run also fitted a power-spectrum slope
$\alpha$ at each coordinate; the distribution is bimodal
(Figure~\ref{fig:grace_run}, bottom centre), with a mean near 0.7 read from the
histogram. The run's per-coordinate table was not saved, so the slope count and the
single-coordinate slope quoted in v3 are withdrawn until the run is repeated. A test
of the prediction needs a method that keeps a steady offset (for example, hotspot
against control coordinates) and an instrument about $10^{6}$ times more sensitive.

\paragraph{Interpretation.}
The GRACE run neither supports nor tests the framework: the predicted
$6.3\times10^{-15}~\mathrm{m\,s^{-2}}$ lies $\sim 6\times10^{5}$ times below the scale
Eq.~\eqref{eq:grace_ul}. The GWOSC sensitivity floor Eq.~\eqref{eq:gwosc_ul} lies $\sim 39$ times
\emph{below} the prediction; it is consistent with the framework only because it
bounds a $\sim 100$~Hz component that the quasi-static prediction does not have. If
a later version of the bulk-leakage model gives $\Delta g$ such a component at
that amplitude, data at this sensitivity could test it. They establish a quantitative
sensitivity floor for next-generation instruments (Stanford quantum
gradiometer at $\leq 10^{-17}~\mathrm{m\,s^{-2}}$, MAGIS-100 commissioning
2027+, GRACE-C 2028+, ESA NGGM 2030s) and define the figure of merit at which
the framework becomes unambiguously falsifiable. Both analysis scripts
(\texttt{TMRB\_GWOSC\_torsion\_search.py} and the GRACE script used for the run,
operating on the identical 45-coordinate set), the GWOSC result file and the GRACE
run's output figure are included in the Zenodo record for this version.

\begin{figure}[ht!]
\centering
\includegraphics[width=0.72\textwidth]{fig5_constraints.pdf}
\caption{Current and projected sensitivity for the predicted
$\Delta g_{\mathrm{TMRB}}=6.3\times10^{-15}~\mathrm{m\,s^{-2}}$ bulk-gravity anomaly:
log-scale sensitivity of operational, planned and proposed missions; the gold band
marks the Gaia-calibrated TMRB prediction; dotted lines mark the GWOSC sensitivity
floor Eq.~\eqref{eq:gwosc_ul} (applies only to a $\sim 100$~Hz component) and the
GRACE scale Eq.~\eqref{eq:grace_ul} (gravity of 1~cm of water). The v3 right panel, a
$z$-score histogram, was produced by a seeded random draw rather than taken from the
run, and is removed.}
\label{fig:constraints}
\end{figure}

\begin{figure}[ht!]
\centering
\includegraphics[width=\textwidth]{fig6_grace_run_2026-04-26.png}
\caption{Output of the 2026-04-26 GRACE run on the real mascon product, reproduced
unchanged with the run script's own panel titles. Each coordinate's series has its
fitted mean, trend and seasonal cycles removed before its mean anomaly and $z$-score are
computed, so both are zero by construction: the $z$-scores cluster at 0 (top right) and
the ring means are 0 (bottom left; dashed line: the prediction). Bottom centre:
power-spectrum slopes $\alpha$. Bottom right: the residual series at one coordinate,
varying at the centimetre level.}
\label{fig:grace_run}
\end{figure}

% ─────────────────────────────────────────────────────────────────────────────
\needspace{6\baselineskip}
\section{Conclusion}
\label{sec:conclusion}

The TMRB framework proposes a resolution of the black-hole information paradox
through three interconnected mechanisms.  First, the $\mathbb{Z}_9$ torsion structure of the
Schoen gyroid compactification provides nine soft-hair charge sectors per torsion
cell on the event horizon — one vacuum ($n=0$) and eight charged ($n=1,\ldots,8$)
— with total Hilbert space dimension $9^{N_{\mathrm{cell}}} = e^{S_{\mathrm{BH}}}$
(Eq.~\eqref{eq:N_hair}).  Second,
the modular-covariant Gaussian weights $w_n = C\exp(-\pi d(n)^2/9)$ generate, beyond the
one-loop logarithm, a universal size-independent torsion correction
$\Delta S_{\mathrm{tors}} = 9\ln 9 + \ln 3 \approx 20.874$
to the Bekenstein--Hawking entropy, where $9\ln 9 = \ln(9^9)$ is the nine-node
hair-lattice log-degeneracy and $\ln 3 = \ln Z_{\mathrm{tors}}$ is the Euclidean
boundary contribution.  Third, a proposed golden-ratio factor $\varphi^{-N_*}$ would
exponentially suppress logarithmic soft-graviton loop corrections (Gap~15), while
holographic consistency is argued from torus partition function factorisation
$Z(\tau)=|\chi_{\mathbb{Z}_9}(\tau)|^2$.  Information-carrying remnants are excluded by a capacity bound; firewall
avoidance and the full unitary map remain open (Gaps~10 and~11).

These results are not isolated: they are the direct UV extension of the same
$\mathbb{Z}_9$ mathematics proposed to close the infrared triangle at $\mathscr{I}^+$
~\cite{BresslerCH2026,BresslerMellin2026}, sets the 230~Myr galactic disk-wave
period~\cite{BresslerGP2026}, and predicts the $\mathbb{Z}_9$ parafermion
spectrum~\cite{BresslerCH2026}.  Section~\ref{sec:fluid_gravity} closes the
UV-to-IR loop through a structural identity: the same parameter $\kappa_{\rm CE}$
that governs the celestial conformal-weight comb at $\mathscr{I}^+$ also fixes
the horizon QNM sideband spacing.  The resulting nine sidebands are too close to
resolve (0.094 line widths at any mass). Their predicted signature, a
mass-independent shortening of the ringdown damping time by about 10--18\% for typical
remnants (Table~\ref{tab:ringdown_spin}), was tested by the X24si method against
GWTC-5.0 and is excluded under the assumptions of Gaps~17 and~18
(Section~\ref{sec:x24si_test}).  The GRACE scale Eq.~\eqref{eq:grace_ul} lies about $10^{6}$ times above the
predicted bulk-gravity-leakage amplitude at the 45 hotspot
coordinates~\cite{BresslerBGL2026}, so GRACE cannot test it; the GWOSC sensitivity floor Eq.~\eqref{eq:gwosc_ul} applies to
a $\sim 100$~Hz component of it, and together they define the sensitivity floor for
next-generation instruments.  The entire TMRB series rests on a single
calibration anchor, $\eta_{\mathrm{diff}} = 0.82\pm0.02$ from Gaia
DR3~\cite{GaiaDR3}.

The universe is a living 9/4-time symphony.  If Gaps~10 and~11 close, with
information conserved at the horizon, the score is complete.

% ─────────────────────────────────────────────────────────────────────────────
\appendix

\needspace{6\baselineskip}
\section{Explicit Microstate Counting}
\label{app:microstate}

The Planck-area decomposition of the horizon gives
$N_{\mathrm{tiles}} = A/(4\ell_{\mathrm{Pl}}^2)$ independent cells.  Each cell
supports a $\mathbb{Z}_9$-valued torsion state $s_i \in \{0,1,\ldots,8\}$, so the
total number of microstates is $9^{N_{\mathrm{tiles}}}$.  The microstate entropy is
\begin{align}
  S_{\mathrm{micro}} &= \ln\!\left(9^{A/(4\ell_{\mathrm{Pl}}^2)}\right)
    = \frac{A}{4\ell_{\mathrm{Pl}}^2}\ln 9. \label{eq:S_micro}
\end{align}
This matches $S_{\mathrm{BH}} = A/(4G_4)$ only at the identification
$G_4 = \ell_{\mathrm{Pl}}^2/\ln 9$, which is not a unit convention but Gap~8.
This is the v1/v2 count; from v3 on the paper tiles the horizon by
$N_{\mathrm{cell}} = A/(4G_4\ln 9)$ cells, giving $S_{\mathrm{micro}} = A/(4G_4)$
directly (Eq.~\eqref{eq:N_hair}).  The
size-independent torsion correction $\Delta S_{\mathrm{tors}} = 9\ln 9 + \ln 3$
arises from two distinct contributions: the nine-node hair-lattice log-degeneracy
$\mathcal{H}_{\mathrm{hair}} = \ln(9^9) = 9\ln 9$ (a Bekenstein-style counting
term for the full $\mathbb{Z}_9$ orbit) and the Euclidean boundary term
$\ln Z_{\mathrm{tors}} = \ln 3$ from the torsion partition function.  The
sign is positive — opposite to the $-\ln\mathcal{D}$ term in
Chern--Simons theory~\cite{LevinWen2006} — because torsion hair adds microstates
rather than removing entanglement degeneracy.

% ── Bibliography ─────────────────────────────────────────────────────────────
\clearpage
\sloppy\raggedright
\begin{thebibliography}{99}

\bibitem{Damour1978}
  T.~Damour,
  \textit{Black-hole eddy currents},
  Phys.\ Rev.\ D \textbf{18}, 3598 (1978).
  \href{https://doi.org/10.1103/PhysRevD.18.3598}{doi:10.1103/PhysRevD.18.3598}

\bibitem{KSS2005}
  P.~K.~Kovtun, D.~T.~Son, A.~O.~Starinets,
  \textit{Viscosity in strongly interacting quantum field theories from black hole physics},
  Phys.\ Rev.\ Lett.\ \textbf{94}, 111601 (2005).
  \href{https://doi.org/10.1103/PhysRevLett.94.111601}{doi:10.1103/PhysRevLett.94.111601}

\bibitem{Hawking1975}
  S.~W.~Hawking,
  \textit{Particle creation by black holes},
  Commun.\ Math.\ Phys.\ \textbf{43}, 199 (1975).
  \href{https://doi.org/10.1007/BF02345020}{doi:10.1007/BF02345020}

\bibitem{AharonovEtal1987}
  Y.~Aharonov, A.~Casher, S.~Nussinov,
  \textit{The unitarity puzzle and Planck mass stable particles},
  Phys.\ Lett.\ B \textbf{191}, 51 (1987).
  \href{https://doi.org/10.1016/0370-2693(87)91320-7}{doi:10.1016/0370-2693(87)91320-7}

\bibitem{Susskind1993}
  L.~Susskind, L.~Thorlacius, J.~Uglum,
  \textit{The stretched horizon and black hole complementarity},
  Phys.\ Rev.\ D \textbf{48}, 3743 (1993).
  \href{https://doi.org/10.1103/PhysRevD.48.3743}{doi:10.1103/PhysRevD.48.3743}

\bibitem{CKW2000}
  V.~Coffman, J.~Kundu, W.~K.~Wootters,
  \textit{Distributed entanglement},
  Phys.\ Rev.\ A \textbf{61}, 052306 (2000).

\bibitem{AMPS2013}
  A.~Almheiri, D.~Marolf, J.~Polchinski, J.~Sully,
  \textit{Black holes: complementarity or firewalls?},
  JHEP \textbf{2013}, 062 (2013).
  \href{https://doi.org/10.1007/JHEP02(2013)062}{doi:10.1007/JHEP02(2013)062}

\bibitem{Penington2020}
  G.~Penington,
  \textit{Entanglement wedge reconstruction and the information paradox},
  JHEP \textbf{2020}, 002 (2020).
  \href{https://doi.org/10.1007/JHEP09(2020)002}{doi:10.1007/JHEP09(2020)002}

\bibitem{HawkingPerryStrominger2016}
  S.~W.~Hawking, M.~J.~Perry, A.~Strominger,
  \textit{Soft hair on black holes},
  Phys.\ Rev.\ Lett.\ \textbf{116}, 231301 (2016).
  \href{https://doi.org/10.1103/PhysRevLett.116.231301}{doi:10.1103/PhysRevLett.116.231301}

\bibitem{Weinberg1965}
  S.~Weinberg,
  \textit{Infrared photons and gravitons},
  Phys.\ Rev.\ \textbf{140}, B516 (1965).
  \href{https://doi.org/10.1103/PhysRev.140.B516}{doi:10.1103/PhysRev.140.B516}

\bibitem{BMS1962}
  H.~Bondi, M.~G.~J.~van der Burg, A.~W.~K.~Metzner,
  \textit{Gravitational waves in general relativity. VII.},
  Proc.\ R.\ Soc.\ London A \textbf{269}, 21 (1962).
  \href{https://doi.org/10.1098/rspa.1962.0161}{doi:10.1098/rspa.1962.0161}

\bibitem{Strominger2014}
  A.~Strominger,
  \textit{On BMS invariance of gravitational scattering},
  JHEP \textbf{2014}, 152 (2014).
  \href{https://doi.org/10.1007/JHEP07(2014)152}{doi:10.1007/JHEP07(2014)152}

\bibitem{Christodoulou1991}
  D.~Christodoulou,
  \textit{Nonlinear nature of gravitation and gravitational-wave experiments},
  Phys.\ Rev.\ Lett.\ \textbf{67}, 1486 (1991).
  \href{https://doi.org/10.1103/PhysRevLett.67.1486}{doi:10.1103/PhysRevLett.67.1486}

\bibitem{Donnay2016}
  L.~Donnay, G.~Giribet, H.~A.~Gonzalez, M.~Pino,
  \textit{Supertranslations and superrotations at the black hole horizon},
  Phys.\ Rev.\ Lett.\ \textbf{116}, 091101 (2016).
  \href{https://doi.org/10.1103/PhysRevLett.116.091101}{doi:10.1103/PhysRevLett.116.091101}

\bibitem{GepnerQiu1987}
  D.~Gepner, Z.-a.~Qiu,
  \textit{Modular invariant partition functions for parafermionic field theories},
  Nucl.\ Phys.\ B \textbf{285}, 423 (1987).
  \href{https://doi.org/10.1016/0550-3213(87)90348-8}{doi:10.1016/0550-3213(87)90348-8}

\bibitem{ZamoFateev1987}
  V.~A.~Fateev, A.~B.~Zamolodchikov,
  \textit{Nonlocal (parafermion) currents in two-dimensional conformal quantum field
    theory and self-dual critical points in $Z_N$-symmetric statistical systems},
  Sov.\ Phys.\ JETP \textbf{62}, 215 (1985) [Zh.\ Eksp.\ Teor.\ Fiz.\ \textbf{89}, 380 (1985)].

\bibitem{AJL2004}
  J.~Ambjorn, J.~Jurkiewicz, R.~Loll,
  \textit{Emergence of a 4D world from causal quantum gravity},
  Phys.\ Rev.\ Lett.\ \textbf{93}, 131301 (2004).
  \href{https://doi.org/10.1103/PhysRevLett.93.131301}{doi:10.1103/PhysRevLett.93.131301}

\bibitem{AJL2005}
  J.~Ambjorn, J.~Jurkiewicz, R.~Loll,
  \textit{Spectral dimension of the universe in quantum gravity},
  Phys.\ Rev.\ D \textbf{72}, 064014 (2005).
  \href{https://doi.org/10.1103/PhysRevD.72.064014}{doi:10.1103/PhysRevD.72.064014}

\bibitem{Vafa1986}
  C.~Vafa,
  \textit{Modular invariance and discrete torsion on orbifolds},
  Nucl.\ Phys.\ B \textbf{273}, 592 (1986).
  \href{https://doi.org/10.1016/0550-3213(86)90565-6}{doi:10.1016/0550-3213(86)90565-6}

\bibitem{DijkgraafWitten1990}
  R.~Dijkgraaf, E.~Witten,
  \textit{Topological gauge theories and group cohomology},
  Commun.\ Math.\ Phys.\ \textbf{129}, 393 (1990).
  \href{https://doi.org/10.1007/BF02096988}{doi:10.1007/BF02096988}

\bibitem{Schoen1970}
  A.~H.~Schoen,
  \textit{Infinite periodic minimal surfaces without self-intersections},
  NASA Technical Note D-5541 (1970).
  \href{https://ntrs.nasa.gov/citations/19700020472}{NTRS:19700020472}

\bibitem{GaiaDR3}
  Gaia Collaboration,
  \textit{Gaia Data Release 3: Summary of the content and survey properties},
  A\&A \textbf{674}, A1 (2023).
  \href{https://doi.org/10.1051/0004-6361/202243940}{doi:10.1051/0004-6361/202243940}

\bibitem{GibbonsHawking1977}
  G.~W.~Gibbons, S.~W.~Hawking,
  \textit{Action integrals and partition functions in quantum gravity},
  Phys.\ Rev.\ D \textbf{15}, 2752 (1977).
  \href{https://doi.org/10.1103/PhysRevD.15.2752}{doi:10.1103/PhysRevD.15.2752}

\bibitem{LevinWen2006}
  M.~Levin, X.-G.~Wen,
  \textit{Detecting topological order in a ground state wave function},
  Phys.\ Rev.\ Lett.\ \textbf{96}, 110405 (2006).
  \href{https://doi.org/10.1103/PhysRevLett.96.110405}{doi:10.1103/PhysRevLett.96.110405}

\bibitem{KitaevPreskill2006}
  A.~Kitaev, J.~Preskill,
  \textit{Topological entanglement entropy},
  Phys.\ Rev.\ Lett.\ \textbf{96}, 110404 (2006).
  \href{https://doi.org/10.1103/PhysRevLett.96.110404}{doi:10.1103/PhysRevLett.96.110404}

\bibitem{DiFrancesco1997}
  P.~Di Francesco, P.~Mathieu, D.~S\'en\'echal,
  \textit{Conformal Field Theory},
  Springer, New York (1997).
  \href{https://doi.org/10.1007/978-1-4612-2256-9}{doi:10.1007/978-1-4612-2256-9}

\bibitem{Page1993}
  D.~N.~Page,
  \textit{Information in black hole radiation},
  Phys.\ Rev.\ Lett.\ \textbf{71}, 3743 (1993).
  \href{https://doi.org/10.1103/PhysRevLett.71.3743}{doi:10.1103/PhysRevLett.71.3743}

\bibitem{Sen2012}
  A.~Sen,
  \textit{Logarithmic corrections to Schwarzschild and other non-extremal
    black hole entropy in different dimensions},
  JHEP \textbf{2013}, 156 (2013).
  \href{https://doi.org/10.1007/JHEP04(2013)156}{doi:10.1007/JHEP04(2013)156}
  \href{https://arxiv.org/abs/1205.0971}{[arXiv:1205.0971]}

\bibitem{Leaver1985}
  E.~W.~Leaver,
  \textit{An analytic representation for the quasi-normal modes of Kerr black holes},
  Proc.\ R.\ Soc.\ Lond.\ A \textbf{402}, 285 (1985).
  \href{https://doi.org/10.1098/rspa.1985.0119}{doi:10.1098/rspa.1985.0119}

\bibitem{Stein2019}
  L.~C.~Stein,
  \textit{qnm: A Python package for calculating Kerr quasinormal modes, separation
    constants, and spherical-spheroidal mixing coefficients},
  J.\ Open Source Softw.\ \textbf{4}, 1683 (2019).
  \href{https://doi.org/10.21105/joss.01683}{doi:10.21105/joss.01683}

\bibitem{BertiCardosoWill2006}
  E.~Berti, V.~Cardoso, C.~M.~Will,
  \textit{On gravitational-wave spectroscopy of massive black holes with the space
    interferometer LISA},
  Phys.\ Rev.\ D \textbf{73}, 064030 (2006).
  \href{https://doi.org/10.1103/PhysRevD.73.064030}{doi:10.1103/PhysRevD.73.064030}

\bibitem{GWOSC2021}
  R.~Abbott et al.\ (LIGO/Virgo/KAGRA Collaborations),
  \textit{Open data from the third observing run of LIGO, Virgo, KAGRA and GEO},
  ApJS \textbf{267}, 29 (2023).
  \href{https://doi.org/10.3847/1538-4365/acdc9f}{doi:10.3847/1538-4365/acdc9f}
  \href{https://gwosc.org}{[gwosc.org]}

\bibitem{GRACE_RL06}
  D.~N.~Wiese, F.~W.~Landerer, M.~M.~Watkins,
  \textit{Quantifying and reducing leakage errors in the JPL RL06 GRACE mascon solution},
  Water Resour.\ Res.\ \textbf{52}, 7490 (2016).
  \href{https://doi.org/10.1002/2016WR019344}{doi:10.1002/2016WR019344}
  ~~~Data product:
  \href{https://podaac.jpl.nasa.gov/dataset/TELLUS_GRAC-GRFO_MASCON_CRI_GRID_RL06.3_V4}{GRCTellus.JPL.200204\_202602.GLO.RL06.3M.MSCNv04CRI.nc}

\bibitem{BresslerUF2026}
  G.~Bressler,
  \textit{TMRB Unified Framework}, v3,
  Zenodo (2026).
  \href{https://doi.org/10.5281/zenodo.19192990}{doi:10.5281/zenodo.19192990}

\bibitem{BresslerCH2026}
  G.~Bressler,
  \textit{Resonant Celestial Holography}, v4,
  Zenodo (2026).
  \href{https://doi.org/10.5281/zenodo.19696857}{doi:10.5281/zenodo.19696857}

\bibitem{BresslerBGL2026}
  G.~Bressler,
  \textit{M-Theory Bulk Gravity Leakage: $\mathbb{Z}_9$ Torsion Nodes, KK Hierarchy,
    and Sub-Femto-$g$ Anomaly Predictions from 11D TMRB}, v10a,
  Zenodo (2026).
  \href{https://doi.org/10.5281/zenodo.19581289}{doi:10.5281/zenodo.19581289}

\bibitem{BresslerGP2026}
  G.~Bressler,
  \textit{The Galactic Piano Is Still Ringing}, v4,
  Zenodo (2026).
  \href{https://doi.org/10.5281/zenodo.19422604}{doi:10.5281/zenodo.19422604}

\bibitem{BresslerQG2026}
  G.~Bressler,
  \textit{TMRB Predictions for NASA QGGPf}, v4,
  Zenodo (2026).
  \href{https://doi.org/10.5281/zenodo.19359426}{doi:10.5281/zenodo.19359426}

\bibitem{BresslerMellin2026}
  G.~Bressler,
  \textit{TMRB Mellin Transform Notebook},
  Zenodo (2026).
  \href{https://doi.org/10.5281/zenodo.19664879}{doi:10.5281/zenodo.19664879}

\bibitem{BresslerPrismatic2026}
  G.~Bressler,
  \textit{Einstein Equations from Prismatic Curvature},
  Zenodo (2026).
  \href{https://doi.org/10.5281/zenodo.20146627}{doi:10.5281/zenodo.20146627}

\bibitem{BresslerX24si2026}
  G.~Bressler,
  \textit{X24si},
  \href{https://x24si.com}{https://x24si.com} (2026).

\bibitem{GWTC5TGR}
  A.~G.~Abac et al.\ (LIGO Scientific, Virgo and KAGRA Collaborations),
  \textit{GWTC-5.0: Tests of General Relativity},
  arXiv:2607.19293 (2026).
  \href{https://arxiv.org/abs/2607.19293}{arXiv:2607.19293}

\bibitem{BresslerSurvey2026}
  G.~Bressler,
  \textit{PTRH Survey/Synthesis: Cross-Citation Map, Master Parameter Table,
    and TRL-3/4 Test Protocol for the Prismatic Torsional Recursive Holographic
    Series},
  Zenodo (2026).
  \href{https://doi.org/10.5281/zenodo.20315669}{doi:10.5281/zenodo.20315669}

\end{thebibliography}

\end{document}
